\documentclass[11pt]{article}
% Shared design for the two minimal papers.
% It follows the layout of the guide "The Riemann Hypothesis and its generalizations"
% (KNT with Fable, 9 November 2026): a designed title block, a short story list at the start,
% coloured headings, boxed statements, idea boxes, and annotated figures.
%
% A paper sets \papertitle, \papersubtitle and \paperrunning before \begin{document}
% and calls \maketitleblock after it.

\usepackage[letterpaper,left=1.12in,right=1.12in,top=1.05in,bottom=1.1in]{geometry}
\usepackage{amsmath,mathtools,amssymb,amsthm}
\usepackage{newpxtext}
\usepackage[varbb]{newpxmath}
\usepackage{fontspec}
\setsansfont{texgyreheros}[Extension=.otf, UprightFont=*-regular, BoldFont=*-bold,
  ItalicFont=*-italic, BoldItalicFont=*-bolditalic, Scale=0.88]
\usepackage{microtype}
\usepackage{parskip}
\usepackage{xcolor}
\usepackage{booktabs,array}
\usepackage{enumitem}
\usepackage{multicol}
\usepackage[titles]{tocloft}
\usepackage{tikz}
\usetikzlibrary{decorations.markings,arrows.meta,positioning,calc}
\usepackage[most]{tcolorbox}
\usepackage{fancyhdr}
\usepackage{titlesec}
\usepackage[font=small,labelfont={bf,sf},labelsep=period,skip=6pt]{caption}
\usepackage{url}

\definecolor{accent}{HTML}{1B4F72}
\definecolor{warm}{HTML}{C0562A}
\definecolor{slate}{HTML}{56636E}
\definecolor{ink}{HTML}{1E1E1E}

\usepackage[colorlinks=true,linkcolor=accent,citecolor=accent,urlcolor=accent]{hyperref}

% Lean identifiers: typewriter, breakable at dots and underscores.
\DeclareUrlCommand\lean{\urlstyle{tt}}

% ---------- layout ----------
\linespread{1.06}
\setlength{\parskip}{0.55em plus 0.1em minus 0.05em}
\setlist{itemsep=0.25em, topsep=0.2em, parsep=0.1em}
\titleformat{\section}{\sffamily\bfseries\Large\color{accent}}{\thesection}{0.7em}{}
\titleformat{\subsection}{\sffamily\bfseries\normalsize\color{ink}}{\thesubsection}{0.6em}{}
\titleformat{\paragraph}[runin]{\sffamily\bfseries\small\color{accent}}{}{0pt}{}[\enspace]
\titlespacing*{\section}{0pt}{22pt}{5pt}
\titlespacing*{\subsection}{0pt}{13pt}{1pt}
\titlespacing*{\paragraph}{0pt}{9pt}{0pt}
\pagestyle{fancy}
\fancyhf{}
\renewcommand{\headrulewidth}{0pt}
\fancyhead[L]{\sffamily\footnotesize\color{slate}\paperrunning}
\fancyhead[R]{\sffamily\footnotesize\color{slate}\thepage}
\setlength{\headheight}{14pt}
\setcounter{tocdepth}{1}
\setlength{\cftbeforesecskip}{2pt}
\renewcommand{\cftsecfont}{\small}
\renewcommand{\cftsecpagefont}{\small\color{slate}}
\renewcommand{\cftsecleader}{\cftdotfill{\cftdotsep}}
\setlength{\cftsecnumwidth}{1.7em}
\setlength{\emergencystretch}{2em}
\numberwithin{equation}{section}

% ---------- title block ----------
\newcommand{\papertitle}{}
\newcommand{\papersubtitle}{}
\newcommand{\paperrunning}{}
\newcommand{\paperauthors}{KNT\enspace\textperiodcentered\enspace Astra\enspace\textperiodcentered\enspace Fable~5.2}
\newcommand{\paperdate}{Draft of 11 November 2026}
\newcommand{\maketitleblock}{%
  \thispagestyle{plain}%
  {\sffamily\bfseries\color{accent}\fontsize{25}{29}\selectfont\raggedright\papertitle\par}%
  \vspace{5pt}%
  {\sffamily\Large\color{slate}\raggedright\papersubtitle\par}%
  \vspace{8pt}%
  {\color{accent}\hrule height 1.2pt}%
  \vspace{7pt}%
  {\sffamily\small\color{ink}\paperauthors\hfill\color{slate}\paperdate\par}%
  \vspace{10pt}}

% Small headings inside the front matter.
\newcommand{\fronthead}[1]{\par\vspace{6pt}{\sffamily\bfseries\normalsize\color{ink}#1}\par\vspace{-2pt}}
\makeatletter
\newcommand{\contentsblock}{%
  \par\noindent\begin{minipage}{\linewidth}%
  \fronthead{Contents}%
  \setlength{\parskip}{0pt}\setlength{\columnsep}{2.2em}\raggedright
  \begin{multicols}{2}
  \@starttoc{toc}
  \end{multicols}
  \end{minipage}\par}
\makeatother

% ---------- statements ----------
% Theorems, lemmas, propositions and corollaries share one counter.
% Each statement is in a tinted box with an accent rule at the left.
\newtheoremstyle{boxed}{0pt}{0pt}{\normalfont}{0pt}{\sffamily\bfseries\small\color{accent}}{.}{0.6em}{}
\theoremstyle{boxed}
\newtheorem{theorem}{Theorem}[section]
\newtheorem{lemma}[theorem]{Lemma}
\newtheorem{proposition}[theorem]{Proposition}
\newtheorem{corollary}[theorem]{Corollary}

\tcbset{guidebox/.style={enhanced, parbox=false, boxrule=0pt, frame hidden,
  arc=1.5pt, left=11pt, right=10pt, top=7pt, bottom=7pt, before skip=11pt, after skip=9pt}}
% Not breakable: a label at the start of a box that breaks at the top of a page can be lost.
\tcbset{statementbox/.style={guidebox, colback=accent!6,
  borderline west={3pt}{0pt}{accent}}}
\tcolorboxenvironment{theorem}{statementbox}
\tcolorboxenvironment{lemma}{statementbox}
\tcolorboxenvironment{proposition}{statementbox}
\tcolorboxenvironment{corollary}{statementbox}

% The main result of a paper: a framed box with a title bar.
% \begin{maintheorem}{name} uses the counter of the theorems, and \label works after it.
\newtcolorbox{mainresult}[1]{enhanced, parbox=false, colback=white, colframe=accent, boxrule=0.9pt,
  arc=1.5pt, left=11pt, right=11pt, top=8pt, bottom=8pt, title={#1},
  fonttitle=\sffamily\bfseries\small, colbacktitle=accent, coltitle=white,
  before skip=12pt, after skip=12pt}
\newenvironment{maintheorem}[1]%
  {\refstepcounter{theorem}\begin{mainresult}{Theorem~\thetheorem\ (#1)}}%
  {\end{mainresult}}

% ---------- proofs ----------
% A proof has a thin rule at the left, so that the eye finds the statements.
\renewcommand{\proofname}{\upshape\sffamily\bfseries\small\color{slate}Proof}
\renewcommand{\qedsymbol}{\textcolor{slate}{\rule{0.55em}{0.55em}}}
\tcolorboxenvironment{proof}{enhanced, breakable, parbox=false, boxrule=0pt, frame hidden,
  colback=white, borderline west={1.2pt}{0pt}{slate!35}, left=11pt, right=0pt, top=1pt, bottom=2pt,
  before skip=6pt, after skip=11pt}
% A step or a case inside a proof.
\newcommand{\step}[1]{{\sffamily\bfseries\small\color{slate}#1}}

% ---------- idea, notation and fine-print boxes ----------
\newtcolorbox{idea}[1]{guidebox, breakable, colback=warm!6, borderline west={3pt}{0pt}{warm!80},
  before upper={{\sffamily\bfseries\footnotesize\color{warm!75!black}IDEA\enspace\textbar\enspace}{\sffamily\bfseries\small #1}\par}}
\newtcolorbox{notation}[1]{guidebox, breakable, colback=black!4, borderline west={3pt}{0pt}{slate!60},
  fontupper=\small,
  before upper={{\sffamily\bfseries\footnotesize\color{slate}NOTATION\enspace\textbar\enspace}{\sffamily\bfseries\small #1}\par}}
\newtcolorbox{fineprint}{guidebox, breakable, colback=white, borderline west={1.5pt}{0pt}{black!25},
  fontupper=\footnotesize\color{black!78}, left=9pt, top=1pt, bottom=1pt,
  before upper={{\sffamily\bfseries Fine print.}\enspace}}

% ---------- figures ----------
\tikzset{
  vtx/.style={circle, fill=ink, inner sep=1.35pt},
  port/.style={circle, draw=accent, fill=white, line width=0.8pt, inner sep=1.5pt},
  edge/.style={line width=0.55pt, draw=ink},
  hot/.style={line width=1.1pt, draw=warm},
  cool/.style={line width=1.1pt, draw=accent},
  region/.style={draw=accent, fill=accent!7, line width=0.6pt, rounded corners=6pt},
  note/.style={font=\sffamily\footnotesize, text=slate},
}

% ---------- macros of the mathematics ----------
\newcommand{\Z}{\mathcal{Z}}
\newcommand{\PP}{\mathbb{P}}
\newcommand{\EE}{\mathbb{E}}
\newcommand{\comp}{\operatorname{comp}}
\newcommand{\tr}{\operatorname{tr}}
\newcommand{\lgk}[1]{\log^{(#1)}}
\newcommand{\lgs}{\log^{*}}
\newcommand{\sm}{\setminus}
\newcommand{\half}{\tfrac12}
\newcommand{\cO}{\mathcal{O}}


% \needroom{length}: start a new page if less than this length is left on the page.
% The pages of this design have no stretch, so the macro measures the room that is left.
\makeatletter
\newcommand{\needroom}[1]{\par\penalty-100\begingroup
  \setlength{\dimen@}{#1}%
  \dimen@ii\pagegoal \advance\dimen@ii-\pagetotal
  \ifdim\dimen@>\dimen@ii \ifdim\dimen@ii>\z@ \vfil\break \fi\fi
  \endgroup}
\makeatother

\renewcommand{\papertitle}{The maximum number of cycles\\ in a graph of large rank}
\renewcommand{\papersubtitle}{The formulas of Horak and Kim are exact for every large rank}
\renewcommand{\paperdate}{Draft of 12 November 2026}
\renewcommand{\paperrunning}{The maximum number of cycles in a graph of large rank}
\hypersetup{pdftitle={The maximum number of cycles in a graph of large rank},pdfauthor={KNT, Astra, Fable 5.2}}

\begin{document}
\maketitleblock

Let $\psi(r)$ be the largest number of cycles in a graph of cyclomatic number $r$. We prove that,
for all sufficiently large $r$,
\[
  \psi(2q+1)=2^{2q}+\bigl(2q+\tfrac12\bigr)2^{q}-3q,\qquad
  \psi(2q+2)=2^{2q+1}+\bigl(2q+\tfrac92\bigr)2^{q}-3q-4 .
\]
Cubic graphs of Horak and Kim have these numbers of cycles. In particular
$\psi(r)\sim2^{r-1}$, which Entringer and Slater conjectured in 1981. The proof uses the
forest estimate of the companion paper~\cite{Forest}. Astra and Fable~5.2 are language models. Section~\ref{sec:status} states the part of each author. The theorem has a
formal proof in Lean~4, and the Lean kernel has checked it.

\fronthead{The proof in seven sentences}
\begin{enumerate}[leftmargin=1.6em]
\item Subdividing an edge keeps the number of cycles, and splitting a vertex of degree four or
more does not lower it. So some cubic multigraph $G$ of rank $r$ has $\psi(r)$ cycles.
\item Fix a small connected set $P$ of vertices. For a cycle that meets $P$, the edges without
an end in $P$ form a forest. By the \emph{forest estimate} of~\cite{Forest}, at most
$\half+o(1)$ of all even subgraphs have this property, and few cycles avoid $P$. So
$\psi(r)\le(\half+o(1))2^{r}$.
\item A cut with at most three edges, or a cycle of length four, divides $G$ into two parts,
and a cycle that meets both parts consists of paths in the parts. A part with more than half of
the possible paths between two ends, by a fixed margin, is a ladder up to a bounded error, and
it has few cycles.
\item So a graph with such a cut or such a cycle is a ring of copies of $K_{2,2}$ and of single
edges, or it has at most $2^{r-1}+O(2^{r/2})$ cycles.
\item A ring has at most $H(r)$ cycles, by a direct count; $H(r)$ is the right side of the
formulas above.
\item In the other case $G$ has girth at least five, and it has fewer than $2^{r-1}$ cycles:
by a second-moment argument, $G$ contains two large pieces with small boundaries, and each piece
has a forest proportion of at most $\frac35$.
\item $H(r)-2^{r-1}$ is at least $(r-2)2^{r/2-1}$. So $\psi(r)=H(r)$ for large $r$.
\end{enumerate}

\contentsblock

\section{The result and the plan}\label{sec:intro}

\subsection{The result}

Throughout, $\log$ is the binary logarithm, and $\lgk{j}$ is its $j$-fold iterate. For a graph
$G$, an \emph{even subgraph} is a set of edges that meets every vertex an even number of times.
With symmetric difference as addition, the even subgraphs form a vector space $\Z(G)$ over
$\mathbb{F}_2$. Its dimension is
\[
  \beta(G)=|E(G)|-|V(G)|+\comp(G),
\]
where $\comp(G)$ is the number of components. We call $\beta(G)$ the \emph{rank} of $G$; other
names are cyclomatic number and cycle rank. A \emph{cycle} is a connected $2$-regular subgraph. We
identify a subgraph without isolated vertices with its edge set. Then a cycle is a nonzero element
of $\Z(G)$, so a graph of rank $r$ has at most $2^{r}-1$ cycles.

Let $\psi(r)$ be the largest number of cycles in a graph of rank $r$.
\begin{itemize}[leftmargin=1.4em]
\item Entringer and Slater~\cite{ES} proved that $2^{r-1}<\psi(r)<2^{r}$ for $r\ge3$, and they
conjectured that $\psi(r)\sim2^{r-1}$.
\item Aldred and Thomassen~\cite{AT} proved that $\psi(r)\le\frac{15}{16}2^{r}$.
\item Alahmadi, Aldred, dela Cruz, Sol\'e and Thomassen~\cite{AACST} proved that a connected
graph with $p$ vertices and more than $2p+3\log(3p)$ edges has at most $2^{r-1}$ cycles.
\item Horak and Kim~\cite{HK} found cubic graphs of rank $r$ with $H(r)$ cycles, where
\begin{equation}\label{eq:H}
  H(2q+1)=2^{2q}+\bigl(2q+\tfrac12\bigr)2^{q}-3q,\qquad
  H(2q+2)=2^{2q+1}+\bigl(2q+\tfrac92\bigr)2^{q}-3q-4 .
\end{equation}
\end{itemize}
For odd rank $r=2q+1$ the graph of Horak and Kim is a ring of $q$ copies of $K_{2,2}$, each joined
to the next by two edges. For even rank, one more edge is inserted across the ring
(Figure~\ref{fig:ring}). We call this graph $R_r$.

\begin{figure}[ht]
\centering
\begin{tikzpicture}
  \def\ri{1.15}\def\ro{1.95}\def\dl{15}
  \foreach \i/\th in {1/18,2/-54,3/-126,4/162}{
    \coordinate (li\i) at (\th+\dl:\ri);
    \coordinate (lo\i) at (\th+\dl:\ro);
    \coordinate (ri\i) at (\th-\dl:\ri);
    \coordinate (ro\i) at (\th-\dl:\ro);
    \draw[cool] (li\i)--(ri\i) (lo\i)--(ro\i) (li\i)--(ro\i) (lo\i)--(ri\i);
  }
  \coordinate (xi) at (90:\ri);
  \coordinate (xo) at (90:\ro);
  \draw[hot] (xi)--(xo);
  \draw[edge] (ri4)--(xi) (ro4)--(xo) (xi)--(li1) (xo)--(lo1);
  \draw[edge] (ri1)--(li2) (ro1)--(lo2) (ri2)--(li3) (ro2)--(lo3) (ri3)--(li4) (ro3)--(lo4);
  \foreach \i in {1,2,3,4}{\foreach \n in {li,lo,ri,ro}{\node[vtx] at (\n\i) {};}}
  \node[vtx] at (xi) {};
  \node[vtx] at (xo) {};
  \node[note, text=warm, anchor=west] at (2.75,1.75) {the further edge (even rank only)};
  \draw[warm, line width=0.4pt] (2.7,1.75) -- (0.12,1.6);
  \node[note, text=accent, anchor=west] at (2.75,-0.55) {a copy of $K_{2,2}$};
  \draw[accent, line width=0.4pt] (2.7,-0.55) -- (1.7,-0.62);
  \node[note, anchor=west] at (2.75,0.6) {two edges join consecutive copies};
  \draw[slate, line width=0.4pt] (2.7,0.6) -- (1.62,0.98);
\end{tikzpicture}
\caption{The graph $R_{10}$: a ring of four copies of $K_{2,2}$ and one further edge. It has rank
$10$ and $H(10)=696$ cycles.}
\label{fig:ring}
\end{figure}

\begin{maintheorem}{the maximum number of cycles}\label{thm:exact}
There is an integer $R$ such that $\psi(r)=H(r)$ for every $r\ge R$.
\end{maintheorem}

In particular $\psi(r)=2^{r-1}+(1+o(1))\,r\,2^{\lfloor(r-1)/2\rfloor}$, so $\psi(r)\sim2^{r-1}$.
The proof gives no value for $R$.

\begin{fineprint}
The theorem is false for small $r$. For example, take the vertices $u_i,v_i$ with
$i\in\mathbb{Z}/26$ and the edges $u_iu_{i+1}$, $u_iv_i$ and $v_iv_{i+10}$. This generalised
Petersen graph has rank $27$, and a direct count by computer gives $67\,959\,973$ cycles, while
$H(27)=67\,325\,913$.
\end{fineprint}

\subsection{The three cases of the proof}

The number of cycles does not change when we delete a vertex of degree at most one, or replace a
vertex of degree two and its two edges by one edge. A vertex of degree at least four can be split
into vertices of degree three, and no cycle is lost. So it suffices to count the cycles of a
connected multigraph $G$ of rank $r$ in which all vertices have degree three
(Lemma~\ref{lem:cubic}). The graphs $R_r$ show that $\psi(r)\ge H(r)$.

For the upper bound there are three cases (Figure~\ref{fig:plan}). A \emph{cyclic cut} is a set of
edges whose deletion leaves two parts that each contain a cycle. The number $H(r)-2^{r-1}$ is of
order $r\,2^{r/2}$, so the three bounds give the theorem (Section~\ref{sec:proof}).

\begin{figure}[ht]
\centering
\begin{tikzpicture}[
  ask/.style={draw=accent, line width=0.7pt, rounded corners=2pt, fill=white,
    font=\sffamily\footnotesize, text width=4.6cm, align=center, inner sep=4.5pt},
  res/.style={fill=accent!7, rounded corners=2pt, font=\sffamily\footnotesize,
    text width=6.3cm, align=left, inner sep=4.5pt},
  ar/.style={-{Stealth[length=4.5pt]}, line width=0.6pt, draw=slate},
  lab/.style={font=\sffamily\scriptsize, text=slate, inner sep=1.5pt},
  node distance=5.2mm]
  \node[ask] (q2) {Does $G$ have a cyclic cut of at most three edges?};
  \node[ask, below=of q2] (q3) {Does $G$ have a cycle of length four?};
  \node[res, right=13mm of q2] (o2) {at most $2^{r-1}+O(2^{r/2})$ cycles\\ \textcolor{slate}{Proposition~\ref{prop:threecut}}};
  \node[res, right=13mm of q3] (o3) {a ring, with at most $H(r)$ cycles, or fewer than $2^{r-1}$ cycles\\ \textcolor{slate}{Theorem~\ref{thm:fourcut}, Corollary~\ref{cor:ringmax}}};
  \node[res, below=of q3, text width=4.6cm, align=center] (o4) {girth at least five, and fewer than $2^{r-1}$ cycles\\ \textcolor{slate}{Theorem~\ref{thm:half}}};
  \foreach \a/\b in {q2/o2,q3/o3} \draw[ar] (\a) -- node[lab, above] {yes} (\b);
  \foreach \a/\b in {q2/q3,q3/o4} \draw[ar] (\a) -- node[lab, right] {no} (\b);
\end{tikzpicture}
\caption{The three cases for a connected cubic multigraph $G$ of large rank $r$.}
\label{fig:plan}
\end{figure}

\subsection{Outside parts and the forest estimate}

Fix a set $P$ of vertices of $G$. For an even subgraph $\xi$ we call the set of its edges without
an end in $P$ its \emph{outside part} (Figure~\ref{fig:outside}). If a cycle meets $P$, then its
outside part is a proper subset of a cycle, so it is a forest.

\begin{figure}[ht]
\centering
\begin{tikzpicture}[x=1cm,y=1cm]
  % the set P: a hexagon
  \foreach \i in {0,...,5} \coordinate (w\i) at ({60*\i+30}:0.9);
  \draw[cool] (w0)--(w1)--(w2)--(w3)--(w4)--(w5)--cycle;
  % outside vertices
  \coordinate (a) at (2.1,1.05);
  \coordinate (b) at (3.3,0.25);
  \coordinate (c) at (2.1,-1.05);
  \coordinate (d) at (-2.2,0.9);
  \coordinate (e) at (-3.0,-0.1);
  \coordinate (f) at (-2.0,-0.9);
  \draw[edge] (w3)--(d) (w2)--(d) (w3)--(f);
  \draw[edge] (d)--(e)--(f) (a)--(c);
  % the cycle C: w0 - a - b - c - w5 - w0 ; outside part thick
  \draw[hot] (a)--(b)--(c);
  \draw[hot, dash pattern=on 2.2pt off 1.6pt] (w5)--(w0) (w0)--(a) (c)--(w5);
  \foreach \p in {w0,w1,w2,w3,w4,w5,a,b,c,d,e,f} \node[vtx] at (\p) {};
  \node[note, text=accent] at (0,0) {$P$};
  \node[note, text=warm, anchor=west] at (3.45,0.25) {the outside part of the cycle $C$: a path};
\end{tikzpicture}
\caption{A set $P$ of six vertices and a cycle $C$ that meets $P$. The edges of $C$ with an end
in $P$ are dashed. The other edges of $C$ form its outside part, which is a path. An even
subgraph whose outside part contains a cycle is not a cycle that meets $P$.}
\label{fig:outside}
\end{figure}

\begin{idea}{cycles and outside parts}
Let $P$ be a small set of vertices such that $G[P]$ is connected. By the \emph{forest estimate}
of the companion paper~\cite{Forest} (Theorem~\ref{thm:forest}), the outside part is a forest for
at most $\half+o(1)$ of all even subgraphs. So at most $(\half+o(1))2^{r}$ cycles meet $P$. A cycle that avoids $P$ is a cycle of
$G-P$, and if $G$ has no bridge, the rank of $G-P$ is at most $r-1-|P|/2$
(Lemma~\ref{lem:marked}). This gives $\psi(r)\le(\half+o(1))2^{r}$ (Corollary~\ref{cor:cycles}).
\end{idea}

In Cases~1 and~2 a cycle consists of paths in the parts of $G$. By the forest estimate,
a part with more than half of the possible paths between two of its ends, by a fixed margin, is a
ladder up to a bounded error, and it has few cycles (Section~\ref{sec:poles}). In Case~2 we
also put a small tree, a \emph{cap}, in the place of the cycle of length four, and apply the
bound $(\half+o(1))2^{r}$ to the graph that arises (Lemma~\ref{lem:cap}). In Case~3 the
forest estimate is not sufficient, because we need a bound below $\half$ and its error term is
positive. In this case a second-moment argument gives two disjoint pieces of $G$ with small
boundaries, and each piece contributes a factor of at most $\frac35$ to the proportion of
forests (Section~\ref{sec:half}).

\subsection{Authorship and verification}\label{sec:status}

KNT is a person. Astra and Fable~5.2 are the language models GPT-6 Astra and Claude Fable~5.2.
KNT directed the work at each step. KNT obtained the forest estimate in September 2026 with
Astra, for a problem of Erd\H{o}s on pancyclic graphs~\cite{KNT}; its proof is
in~\cite{Forest}. For the number of cycles, Astra found the path structure of
Section~\ref{sec:poles}, the gap theorem, the count for rings, the cap, and the first proof of
the last case. Fable~5.2 found the lemma on layers, the chain of separators and the pieces at
girth five, and wrote this text. The present forms of the proofs contain steps of both models.

Theorem~\ref{thm:exact} has a formal proof in Lean~4. The statement
\lean{Erdos1016.ExactMaximum.exactMaximum_threeCases} says that there is a number $R$ such that
\lean{maxCycles r = horakKim r} for every $r\ge R$. Here \lean{maxCycles r} is the largest number
of cycles of a simple graph with cyclomatic number $r$, and \lean{horakKim r} is $H(r)$. The Lean
kernel accepts the proof with the axioms \lean{propext}, \lean{Classical.choice} and
\lean{Quot.sound} only. The files are in~\cite{Repo}, at commit \texttt{b7f9fc9}.

The formal proof has the three cases of Section~\ref{sec:proof}, and it proves
Theorem~\ref{thm:half} at girth five. Some of its steps are not those of the text: a cut of
size four has two arbitrary sides, Lemma~\ref{lem:fullfive} has a proof without ladders,
and the end of the proof of Theorem~\ref{thm:half} uses short cycles at many scales. So the
kernel check
covers the theorem, and it does not cover the text of this paper. No referee has read the text.

\section{Tools}\label{sec:tools}

\begin{notation}{multigraphs, cuts and regions}
\begin{itemize}[leftmargin=1.2em]
\item From here on a graph may have loops and parallel edges. We say \emph{multigraph} to stress
this, and \emph{simple graph} to exclude it. A loop adds two to the degree of its vertex. A loop is
a cycle of length one, and two parallel edges form a cycle of length two.
\item A multigraph is \emph{cubic} if every vertex has degree three, and \emph{subcubic} if every
vertex has degree at most three. The \emph{girth} is the length of a shortest cycle.
\item $c(G)$ is the number of cycles of $G$, and $c_e(G)$ is the number of cycles through an edge
$e$. If we replace an edge by a path, the rank and the number of cycles do not change. So
$\psi(r)$ is also the largest number of cycles of a multigraph of rank $r$.
\item For a set $A$ of vertices, $G[A]$ is the induced subgraph, $G-A$ is $G[V(G)\sm A]$, and
$\delta A$ is the set of edges with exactly one end in $A$.
\item A \emph{region} is a nonempty set $A$ such that $G[A]$ is connected. A region is
\emph{cyclic} if $G[A]$ contains a cycle.
\item A \emph{cyclic cut} is a set $\delta X$ such that $G[X]$ and $G-X$ each contain a cycle; its
\emph{size} is $|\delta X|$.
\item Throughout, $\xi$ is a uniformly random element of $\Z(G)$.
\end{itemize}
\end{notation}

\subsection{Outside parts}\label{sec:count}

Let $W$ be a subgraph of $G$, with $r=\beta(G)$ and $t=\beta(W)$. For an even subgraph
$\xi$ we call $\xi\sm E(W)$ its \emph{outside part}. The map $\xi\mapsto\xi\sm E(W)$ is linear on
$\Z(G)$ with kernel $\Z(W)$. So there are $2^{r-t}$ outside parts, each is the outside part of
exactly $2^{t}$ even subgraphs, and the outside part of a uniform $\xi$ is uniform. Put
\[
  p(G,W)=\PP\bigl(\xi\sm E(W)\text{ is a forest}\bigr),\qquad
  q(G,P)=\PP\bigl(\xi\cap E(G-P)\text{ is a forest}\bigr)
\]
for a subgraph $W$ and a set $P$ of vertices.

\begin{lemma}[Edge recursion]\label{lem:edge}
Let $G$ be a connected multigraph of rank $r$, let $W$ be a cycle of $G$, and let $e$ be an edge of
$W$. Then $G-e$ is connected and has rank $r-1$, and
\[
  c(G)=c(G-e)+c_e(G),\qquad c_e(G)\;\le\;p(G,W)\,2^{r-1}.
\]
\end{lemma}

\begin{proof}
The edge $e$ lies on a cycle, so $G-e$ is connected with one edge less. A cycle of $G$ avoids $e$
or contains $e$. This is the first formula.

Let $C$ and $C'$ be cycles through $e$ with the same outside part. Then $C+C'$ is an even subgraph
inside $W$, so it is empty or equal to $W$. But $C+W$ does not contain $e$. So $C'=C$. Hence the
cycles through $e$ have distinct outside parts. Each of these is a proper subset of a cycle, so
it is a forest. There are $p(G,W)\,2^{r-1}$ outside parts that are forests, because
$\beta(W)=1$.
\end{proof}

\subsection{Cuts and links}\label{sec:links}

For a region $A$ let $\cO_A$ be the event $\xi\cap\delta A=\emptyset$.

\begin{lemma}\label{lem:second}
Let $Z\ge0$ be a random variable with $\mu=\EE Z>0$ and $\EE Z^{2}\le2\mu^{2}+K\mu$. Then
$\PP(Z=0)\le\half+K/(4\mu)$.
\end{lemma}

\begin{proof}
By the Cauchy--Schwarz inequality $\mu^{2}\le\EE Z^{2}\cdot\PP(Z>0)$, so
$\PP(Z=0)\le1-\mu/(2\mu+K)\le\half+K/(4\mu)$.
\end{proof}

Let $A$ and $B$ be disjoint regions of a connected multigraph $G$. A \emph{link} between them is an
edge with one end in $A$ and the other in $B$, or a component of $G-(A\cup B)$ that is joined by an
edge to $A$ and by an edge to $B$ (Figure~\ref{fig:links}). Let $\ell(A,B)$ be the number of links.

\begin{figure}[ht]
\centering
\begin{tikzpicture}
  \draw[region] (0,0) ellipse [x radius=0.75, y radius=1.0];
  \draw[region] (6,0) ellipse [x radius=0.75, y radius=1.0];
  \draw[draw=warm, fill=warm!8, line width=0.6pt] (3,-0.7) ellipse [x radius=0.85, y radius=0.38];
  \draw[draw=slate!70, fill=black!4, line width=0.6pt] (-2.9,0.2) ellipse [x radius=0.85, y radius=0.38];
  \draw[draw=slate!70, fill=black!4, line width=0.6pt] (8.9,-0.2) ellipse [x radius=0.85, y radius=0.38];
  \node[text=accent] at (0,0) {$A$};
  \node[text=accent] at (6,0) {$B$};
  \node[text=warm] at (3,-0.7) {$Y_1$};
  \node[text=slate] at (-2.9,0.2) {$Y_2$};
  \node[text=slate] at (8.9,-0.2) {$Y_3$};
  \draw[hot] (0.45,0.55) -- node[above=-1pt, text=warm] {\small$e$} (5.55,0.55);
  \draw[edge] (0.5,-0.45) -- (2.25,-0.58);
  \draw[edge] (0.4,-0.75) -- (2.3,-0.86);
  \draw[edge] (5.5,-0.5) -- (3.75,-0.63);
  \draw[edge] (-0.5,0.35) -- (-2.12,0.32);
  \draw[edge] (-0.55,0.05) -- (-2.1,0.1);
  \draw[edge] (6.55,-0.15) -- (8.07,-0.2);
  \foreach \p in {(0.45,0.55),(5.55,0.55),(0.5,-0.45),(2.25,-0.58),(0.4,-0.75),(2.3,-0.86),(5.5,-0.5),
                  (3.75,-0.63),(-0.5,0.35),(-2.12,0.32),(-0.55,0.05),(-2.1,0.1),(6.55,-0.15),(8.07,-0.2)}
    \node[vtx] at \p {};
  \node[note, text=warm] at (3,1.05) {two links: the edge $e$ and the component $Y_1$};
  \node[note] at (-2.9,-0.45) {not a link};
  \node[note] at (8.9,-0.85) {not a link};
\end{tikzpicture}
\caption{Links between two regions $A$ and $B$. The edge $e$ and the component $Y_1$ of
$G-(A\cup B)$ are links. The components $Y_2$ and $Y_3$ are joined to one region only, so they are
not links. Here $\ell(A,B)=2$.}
\label{fig:links}
\end{figure}

\begin{lemma}[Link formula]\label{lem:link}
Let $A$ and $B$ be disjoint regions of a connected multigraph $G$. Then
$\PP(\cO_A)=2^{\comp(G-A)-|\delta A|}\ge2^{-|\delta A|}$ and
\[
  \PP(\cO_A\cap\cO_B)=2^{\ell(A,B)-1}\,\PP(\cO_A)\,\PP(\cO_B).
\]
Conditionally on $\cO_A$, the set $\xi\cap E(G[A])$ is uniform on $\Z(G[A])$. Conditionally on
$\cO_A\cap\cO_B$, the sets $\xi\cap E(G[A])$ and $\xi\cap E(G[B])$ are independent and uniform on
$\Z(G[A])$ and $\Z(G[B])$.
\end{lemma}

\begin{proof}
The even subgraphs that avoid $\delta A$ are the sums of an even subgraph of $G[A]$ and an even
subgraph of $G-A$. So $\PP(\cO_A)=2^{\beta(G[A])+\beta(G-A)-\beta(G)}$, and a count of edges and
vertices shows that the exponent is $\comp(G-A)-|\delta A|$. In the same way
\[
  \PP(\cO_A\cap\cO_B)=2^{\beta(G[A])+\beta(G[B])+\beta(G-A-B)-\beta(G)}
     =2^{\comp(G-A-B)+1-|\delta A|-|\delta B|+e(A,B)},
\]
where $e(A,B)$ is the number of edges between $A$ and $B$. Let $a$, $b$ and $k$ be the numbers of
components of $G-A-B$ that are joined to $A$ only, to $B$ only, and to both. Then
$\comp(G-A-B)=a+b+k$, $\comp(G-A)=a+1$ and $\comp(G-B)=b+1$. The formula follows, because
$\ell(A,B)=k+e(A,B)$. The conditional laws are read off from the two decompositions as direct sums.
\end{proof}

\begin{idea}{the correlation of two cut events}
By Lemma~\ref{lem:link}, two cut events are independent when the regions have one link, and
their joint probability is twice the product when the regions have two links. The constant
$\half$ of the forest estimate comes from this factor two, through Lemma~\ref{lem:second}.
\end{idea}

Pairs with three or more links are rare:

\begin{lemma}[Exceptional partners]\label{lem:exc}
Let $A$ be a region of a connected multigraph $G$, and let $B_1,\dots,B_k$ be pairwise disjoint
regions, disjoint from $A$, with $\ell(A,B_i)\ge3$ for all $i$. Then $k\le|\delta A|$.
\end{lemma}

\begin{proof}
Contract $A$ to a vertex $a$ and each $B_i$ to a vertex $b_i$, and delete the loops that arise.
This changes neither $d=|\delta A|$ nor the numbers $\ell(A,B_i)$, because the sets $B_j$ with
$j\ne i$ are connected and are contracted inside components of $G-(A\cup B_i)$. Let $F$ be a
spanning forest of $G-a$, and for every edge $e=au$ at $a$ attach a new leaf $\tau_e$ to $u$. Call
the resulting forest $F^{+}$ and its $d$ new leaves \emph{terminals}.

Fix $i$. Every link between $a$ and $b_i$ yields a path of $F^{+}$ from $b_i$ to a terminal. For an
edge $e=ab_i$ it is the edge $b_i\tau_e$. For a component $Y$ of $G-\{a,b_i\}$ joined to both,
choose an edge $e=au$ with $u\in Y$ and take the path of $F^{+}$ from $b_i$ to $\tau_e$; after its
first vertex it stays in $Y\cup\{\tau_e\}$. Paths that come from different links start with
different edges at $b_i$, and the union of any two of them is a path between two terminals. Let
$F^{\circ}\subseteq F^{+}$ be the union of all paths between terminals. As $\ell(a,b_i)\ge3$, the
vertex $b_i$ has degree at least $3$ in $F^{\circ}$. All leaves of $F^{\circ}$ are terminals, and a
forest with at most $d$ leaves has at most $d-2$ vertices of degree at least $3$. Hence $k\le d-2$
if $k\ge1$.
\end{proof}

\begin{lemma}[Many small cuts]\label{lem:smallcuts}
Let $A_1,\dots,A_k$ be pairwise disjoint cyclic regions of a connected multigraph $G$ with
$|\delta A_i|\le d$ for all $i$. Then
\[
  \PP\bigl(\xi\cap E(G[A_i])\text{ is a forest for every }i\bigr)\;\le\;\frac12+\frac{(d+1)\,2^{d-1}}{k}.
\]
\end{lemma}

\begin{proof}
Let $E_i$ be the event that $\cO_{A_i}$ holds and $\xi\cap E(G[A_i])\ne\emptyset$. On $E_i$ the set
$\xi\cap E(G[A_i])$ is a nonempty even subgraph, hence not a forest. So it suffices to bound
$\PP(Z=0)$ for $Z=\sum_i\mathbf{1}_{E_i}$. By Lemma~\ref{lem:link},
$p_i:=\PP(E_i)=\PP(\cO_{A_i})\bigl(1-2^{-\beta(G[A_i])}\bigr)\ge2^{-d-1}$ and
$\PP(E_i\cap E_j)=2^{\ell(A_i,A_j)-1}p_ip_j$ for $i\ne j$. If $\ell(A_i,A_j)\le2$, this is at most
$2p_ip_j$. Otherwise we use $\PP(E_i\cap E_j)\le p_i$; by Lemma~\ref{lem:exc} this happens for at
most $d$ indices $j$, for each $i$. Hence $\EE Z^{2}\le\mu+2\mu^{2}+d\mu$ with
$\mu=\EE Z\ge k\,2^{-d-1}$, and Lemma~\ref{lem:second} gives the claim.
\end{proof}

\subsection{Localisation}\label{sec:loc}

By the next lemma, the part of $\xi$ inside a set $X$ has the same distribution as the part of a
random even subgraph of a smaller graph outside one cycle.

\begin{lemma}[Localisation]\label{lem:loc}
Let $G$ be a connected multigraph, and let $X$ be a nonempty proper set of vertices such that
$G-X$ is connected. Put $d=|\delta X|$. Form $G_X$ from $G[X]$ as follows. For each edge of
$\delta X$ add a new vertex and join it to the end of that edge in $X$. Join the new vertices by a
cycle $W_X$ through all of them; if $d<3$, first add $3-d$ further new vertices on that cycle. Then
\begin{enumerate}
\item[\textup{(a)}] $G_X$ is connected, and if every vertex of $X$ has degree three in $G$, then
$\beta(G_X)=(|X|+d)/2+1$;
\item[\textup{(b)}] the set $\xi\cap E(G[X])$ has the same distribution for $\xi$ uniform on
$\Z(G)$ and for $\xi$ uniform on $\Z(G_X)$;
\item[\textup{(c)}] an even subgraph $\xi$ of $G_X$ has a forest outside $E(W_X)$ if and only if
$\xi\cap E(G[X])$ is a forest.
\end{enumerate}
\end{lemma}

\begin{proof}
(a) Every component of $G[X]$ has an edge to $G-X$, so the cycle $W_X$ joins all components.
Without the extra vertices, $G_X$ has $d$ more vertices and $2d$ more edges than $G[X]$, so
$\beta(G_X)=|E(G[X])|+2d-|X|-d+1$. Extra vertices on the cycle add as many edges as vertices. If
all degrees in $X$ are three, then $3|X|=2|E(G[X])|+d$.

(b) Contract the connected set $V(G)\sm X$ of $G$ to one vertex, and delete the loops that arise.
Call the result $G/X$. An even subgraph of $G$ restricts to an even subgraph of $G/X$. Every even
subgraph of $G/X$ arises in this way: the missing edges inside the contracted set can be chosen in
a spanning tree of that set. So the restriction is a surjective linear map $\Z(G)\to\Z(G/X)$, and
it sends the uniform distribution to the uniform distribution. Now contract $W_X$ in $G_X$. The
result is the same multigraph $G/X$, with the new edges in the place of the edges of $\delta X$.

(c) The edges of $G_X$ outside $W_X$ are the edges of $G[X]$ and the $d$ new edges at $X$. In
$G_X-E(W_X)$ each of these new edges has an end of degree one, so it lies on no cycle.
\end{proof}


\section{The leading term}\label{sec:leading}

This section proves that $\psi(r)\le(\half+o(1))\,2^{r}$. The later sections use this bound and
the two lemmas.

\begin{lemma}[Cubic multigraphs]\label{lem:cubic}
For every $r\ge2$ there is a connected cubic multigraph of rank $r$ with exactly $\psi(r)$ cycles.
\end{lemma}

\begin{proof}
The numbers of cycles of the graphs of rank $r$ are integers below $2^{r}$, so some graph $G_0$
attains $\psi(r)$. If $G_0$ is not connected, identify a vertex of one component with a vertex of
another. This changes neither the rank nor the set of cycles. Now delete a vertex of degree one, or
replace a vertex of degree two and its two edges by one edge, as long as this is possible. These
steps keep the rank and the number of cycles. The result is a connected multigraph with all
degrees at least three, because $r\ge2$.

Let $v$ be a vertex of degree $d\ge4$. Replace $v$ by a path $v^{1}v^{2}\cdots v^{d-2}$, and
distribute the $d$ edge ends at $v$ along this path: two at $v^{1}$, one at each of
$v^{2},\dots,v^{d-3}$, and two at $v^{d-2}$. The new multigraph has $d-3$ more vertices and $d-3$
more edges, so its rank is $r$. A cycle through $v$ uses two edge ends at $v$. Together with the
part of the path between these two ends it is a cycle of the new multigraph. Different cycles give
different cycles, and the cycles that avoid $v$ remain. So the number of cycles does not decrease,
and it cannot exceed $\psi(r)$. Repeat this for every vertex of degree at least four.
\end{proof}

\begin{lemma}[Cycles and a marked set]\label{lem:marked}
Let $G$ be a connected cubic multigraph of rank $r$ without a bridge. Let $P$ be a set of $p$
vertices, not all vertices of $G$, such that $G[P]$ is connected. Then
\[
  c(G)\;\le\;q(G,P)\,2^{r}+2^{\,r-1-p/2}.
\]
\end{lemma}

\begin{proof}
Let $Z$ be a cycle of $G$ that meets $P$. Its edges in $G-P$ form a proper subset of $Z$, hence a
forest. So at most $q(G,P)\,2^{r}$ cycles meet $P$. The other cycles are cycles of $G-P$, and
there are at most $2^{\beta(G-P)}$ of them. Let $a$ be the number of edges of $G[P]$. Then $3p-2a$
edges leave $P$, and $G-P$ arises from $G$ by the deletion of $p$ vertices and $3p-a$ edges. So
$\beta(G-P)=r-1-2p+a+\comp(G-P)$. The graph $G$ has no bridge, so every component of $G-P$
contains the ends of at least two of the edges that leave $P$. Hence $\comp(G-P)\le(3p-2a)/2$ and
$\beta(G-P)\le r-1-p/2$.
\end{proof}

\begin{corollary}[Leading term]\label{cor:cycles}
There is $r_1$ such that every graph of rank $r\ge r_1$ has at most
$\bigl(\half+21/\lgk3r\bigr)2^{r}$ cycles. Consequently $\psi(r)\sim2^{r-1}$.
\end{corollary}

\begin{proof}
Let $r$ be large, and let $G$ be a connected cubic multigraph of rank $r$ with $\psi(r)$ cycles
(Lemma~\ref{lem:cubic}). Suppose first that $G$ has a bridge $e$, and let $b_1$ and $b_2$ be the
ranks of the two components of $G-e$. A component with $n$ vertices has $(3n-1)/2\ge n$ edges, so
$b_1,b_2\ge1$, and $b_1+b_2=r$. Every cycle lies in one component. So
$c(G)\le(2^{b_1}-1)+(2^{b_2}-1)\le2^{r-1}$.

In the other case choose a set $P$ of $\lceil4\log r\rceil$ vertices such that $G[P]$ is
connected. Then $|P|\le2r/\log^{2}r$. By Theorem~\ref{thm:forest} and Lemma~\ref{lem:marked},
\[
  c(G)\;\le\;\Bigl(\frac12+\frac{20}{\lgk3r}\Bigr)2^{r}+2^{\,r-1-2\log r}
  \;\le\;\Bigl(\frac12+\frac{21}{\lgk3r}\Bigr)2^{r}.
\]
Finally $\psi(r)>2^{r-1}$ for $r\ge3$~\cite{ES}; the rings of Section~\ref{sec:rings} also show this.
\end{proof}

\section{Layers}\label{sec:layers}

The later sections use Theorem~\ref{thm:forest} in two ways: by Corollary~\ref{cor:cycles}, and
by the lemma of this section.

\begin{idea}{the content of the lemma}
If the outside part is a forest for more than half of the even subgraphs, by a fixed margin, then
every distance layer around the marked set is small. In the other case a window of consecutive layers is a large set
with a small boundary, and the forest estimate applies inside this set. We use the lemma for
paths in Section~\ref{sec:poles} and for pieces in Section~\ref{sec:half}.
\end{idea}

Let $G$ be a multigraph and $P$ a set of vertices. A \emph{level function} is a map
$\lambda\colon V(G)\to\mathbb{Z}$ such that $|\lambda(u)-\lambda(w)|\le1$ for every edge $uw$ that
has an end outside $P$. The distance from a fixed vertex is a level function. For
$\delta\in(0,\half)$ let $R_\delta\ge2^{100}$ be an integer such that every $r\ge R_\delta$ is in
the range of Theorem~\ref{thm:forest} and satisfies $20/\lgk3r\le\delta$. For an integer $k\ge1$
put
\[
  \Theta_\delta(k)=\max\bigl\{2R_\delta,\;2^{100}(k+1)^{2}\bigr\}.
\]

\begin{lemma}[Layers]\label{lem:width}
Let $G$ be a connected subcubic multigraph, and let $P$ be a nonempty set of vertices such that
$G[P]$ is connected and every vertex outside $P$ has degree three. Let $\lambda$ be a level
function such that at most $k$ vertices of $P$ have any given level. Let $\delta\in(0,\half)$, and
suppose that $q(G,P)>\half+\delta$. Then, for every $j$, fewer than $\Theta_\delta(k)$ vertices
outside $P$ have level $j$.
\end{lemma}

\begin{figure}[ht]
\centering
\begin{tikzpicture}[x=1cm,y=1cm]
  % the graph G
  \draw[draw=slate!70, fill=black!4, line width=0.6pt] (0,0) ellipse [x radius=1.25, y radius=1.15];
  \draw[draw=slate!70, fill=black!12, line width=0.6pt] (-0.4,0.15) ellipse [x radius=0.45, y radius=0.4];
  \node[text=slate] at (-0.4,0.15) {\small$P$};
  \node[note] at (0.15,-0.75) {$C_0$};
  \draw[region] (4.0,0) ellipse [x radius=1.45, y radius=1.15];
  \draw[accent, line width=0.5pt, dash pattern=on 2pt off 2pt] (4.4,0.1) ellipse [x radius=0.75, y radius=0.6];
  \node[text=accent] at (4.4,0.1) {$X$};
  \node[note, text=accent] at (4.0,-0.85) {$X^{+}$};
  \foreach \y in {0.6,0.2,-0.2,-0.6}{
    \draw[edge] (1.05,\y) -- (2.8,\y);
    \node[vtx] at (1.05,\y) {};
    \node[vtx] at (2.8,\y) {};}
  \node[note] at (1.95,-1.6) {the graph $G$};
  % the multigraph Gamma
  \begin{scope}[xshift=8.3cm]
    \draw[hot] (0.5,0) ellipse [x radius=0.33, y radius=0.85];
    \draw[region] (3.7,0) ellipse [x radius=1.45, y radius=1.15];
    \draw[accent, line width=0.5pt, dash pattern=on 2pt off 2pt] (4.1,0.1) ellipse [x radius=0.75, y radius=0.6];
    \node[text=accent] at (4.1,0.1) {$X$};
    \node[note, text=accent] at (3.7,-0.85) {$X^{+}$};
    \foreach \y/\x in {0.6/0.73,0.2/0.82,-0.2/0.82,-0.6/0.73}{
      \draw[edge] (\x,\y) -- (2.5,\y);
      \node[vtx] at (\x,\y) {};
      \node[vtx] at (2.5,\y) {};}
    \node[note, text=warm] at (0.5,1.15) {$W$};
    \node[note, text=warm] at (0.55,-1.2) {$P^{+}$};
    \node[note] at (1.9,-1.6) {the multigraph $\Gamma$};
  \end{scope}
\end{tikzpicture}
\caption{Step~1 in the proof of Lemma~\ref{lem:width}. In $G$ the set $X^{+}$ is joined by $d^{+}$ edges to the
connected rest $C_0$, which contains $P$. In $\Gamma$ a cycle $W$ through $d^{+}$ new vertices
takes the place of $C_0$. The even subgraphs of $G$ and of $\Gamma$ give the same distribution
of the edges inside $X^{+}$.}
\label{fig:local}
\end{figure}

\begin{proof}
Put $\Phi(m)=m/\log^{2}m$ and $\Theta=\Theta_\delta(k)$.

\step{Step 1: no large set outside $P$ has a small boundary.} We claim that no set
$X\subseteq V(G)\sm P$ has $|X|=m\ge\Theta$ and $|\delta X|\le\Phi(m)/4$. Suppose that $X$ is such a
set. Let $C_0$ be the component of $G-X$ that contains $P$, and put $X^{+}=V(G)\sm V(C_0)$. Then
$X\subseteq X^{+}$, the set $X^{+}$ avoids $P$, and $G-X^{+}=C_0$ is connected. Different components
of $G-X$ have no edge between them, so $\delta X^{+}\subseteq\delta X$. Every vertex of $X^{+}$ has
degree three. Put $d^{+}=|\delta X^{+}|\ge1$.

We form a multigraph $\Gamma$ from $G[X^{+}]$ (Figure~\ref{fig:local}). For each edge of $\delta X^{+}$ add a new vertex
and join it to the end of that edge in $X^{+}$. Join the $d^{+}$ new vertices by a cycle $W$
through all of them, and let $P^{+}$ be the set of the new vertices. Every component of
$G[X^{+}]$ has an edge to $C_0$, so $\Gamma$ is connected. It is subcubic, $\Gamma[P^{+}]$ is
connected, the vertices outside $P^{+}$ have degree three, and $\Gamma-P^{+}=G[X^{+}]$. The
multigraph $\Gamma$ has $d^{+}$ more vertices and $2d^{+}$ more edges than $G[X^{+}]$, and
$3|X^{+}|=2|E(G[X^{+}])|+d^{+}$. So its rank is
\[
  \rho=\frac{|X^{+}|+d^{+}}2+1\;\ge\;\frac m2\;\ge\;R_\delta .
\]
The function $\Phi$ is increasing on $[8,\infty)$, and $\Phi(m)/4\le\Phi(m/2)$. So
$d^{+}\le\Phi(m)/4\le\Phi(m/2)\le\Phi(\rho)$, and Theorem~\ref{thm:forest} with the choice of
$R_\delta$ gives $q(\Gamma,P^{+})\le\half+\delta$.

The set $\xi\cap E(G[X^{+}])$ has the same distribution for $\xi$ uniform on $\Z(G)$ and for
$\xi$ uniform on $\Z(\Gamma)$. Indeed, contract the connected set $V(C_0)$ of $G$ to one vertex,
and delete the loops that arise. Call the result $G'$. An even subgraph of $G$ restricts to an
even subgraph of $G'$. Every even subgraph of $G'$ arises in this way: the missing edges inside
the contracted set can be chosen in a spanning tree of that set. So the restriction is a
surjective linear map $\Z(G)\to\Z(G')$, and it sends the uniform distribution to the uniform
distribution. The contraction of $W$ in $\Gamma$ gives the same multigraph $G'$, with the new
edges in the place of the edges of $\delta X^{+}$. Hence
\[
  \PP\bigl(\xi\cap E(G[X^{+}])\text{ is a forest}\bigr)\;=\;q(\Gamma,P^{+})\;\le\;\frac12+\delta .
\]
But $G[X^{+}]$ is a subgraph of $G-P$, so this probability is at least $q(G,P)>\half+\delta$. This
contradiction proves the claim.

\step{Step 2: a window of levels.} Let $m_j$ be the number of vertices outside $P$ of level $j$.
For $i\le i'$ let $X(i,i')$ be the set of vertices outside $P$ with a level in $[i,i']$. An edge
that leaves $X(i,i')$ ends outside $P$ at level $i-1$ or $i'+1$, or it ends in $P$ at a level in
$[i-1,i'+1]$. So
\begin{equation}\label{eq:window}
  |\delta X(i,i')|\;\le\;3m_{i-1}+3m_{i'+1}+3k(i'-i+3).
\end{equation}
Suppose that $m_j\ge\Theta$ for some $j$. Start with the window $i=i'=j$. As long as a level next
to the window has more than $m^{3/4}$ vertices outside $P$, where $m=|X(i,i')|$ is the present
size, add that level to the window. The graph is finite, so this stops. A step increases
$m^{1/4}$ by more than $\frac16$, because $(t+\frac16)^{4}\le t^{4}+t^{3}$ for $t\ge1$. Let $\mu$
be the final size. Then there were fewer than $6\mu^{1/4}$ steps, and both levels next to the
final window have at most $\mu^{3/4}$ vertices outside $P$. We have $k\le\sqrt\mu$, because
$\mu\ge\Theta\ge2^{100}(k+1)^{2}$. By~\eqref{eq:window},
\[
  |\delta X(i,i')|\;\le\;6\mu^{3/4}+3k\bigl(6\mu^{1/4}+3\bigr)\;\le\;33\,\mu^{3/4}
  \;\le\;\frac{\mu}{4\log^{2}\mu}=\frac{\Phi(\mu)}4 .
\]
The last inequality says that $132\log^{2}\mu\le\mu^{1/4}$, and this holds for
$\mu\ge2^{100}$. So the final window contradicts Step~1.
\end{proof}


\section{Poles with many paths}\label{sec:poles}

A cut of a cubic graph divides it into two parts, and a cycle that crosses the cut is made of
paths in the two parts. So we need to count paths between prescribed ends.

\begin{idea}{the results of this section}
First, if more than half of the possible paths between two ends are present, by a fixed margin,
then the part looks like a ladder, and it has few cycles (Theorem~\ref{thm:fewcycles}). Second,
the number of paths between two ends is the largest possible, or it is smaller by a fixed factor
(Theorem~\ref{thm:gap}).
\end{idea}

\subsection{Poles}

A \emph{pole} is a connected multigraph without loops in which every vertex has degree two or
three. Its vertices of degree two are its \emph{ports}. A \emph{$k$-pole} is a pole with exactly
$k$ ports. By a count of edge ends, a $k$-pole of rank $h$ has $2h+k-2$ vertices.

For two vertices $s\ne t$ of a multigraph $N$ let $P_{st}(N)$ be the number of paths from $s$ to
$t$. An \emph{$\{s,t\}$-join} is a set of edges in which $s$ and $t$ have odd degree and all other
vertices have even degree. If $N$ is connected of rank $b$, there are exactly $2^{b}$ of them: they
exist, and the sum of two of them is even. Every path from $s$ to $t$ is an $\{s,t\}$-join, so
\begin{equation}\label{eq:joinbound}
  P_{st}(N)\;\le\;2^{b}.
\end{equation}
Call the pair $s,t$ \emph{full} in $N$ if $P_{st}(N)=2^{\beta(N)}$.

\begin{lemma}[Joins]\label{lem:joins}
Let $N$ be a pole of rank $b$, and let $s,t$ be two of its ports. Every $\{s,t\}$-join is the union
of a path from $s$ to $t$ and of cycles, and these are pairwise vertex-disjoint. So the pair is
full if and only if every cycle of $N$ meets every path from $s$ to $t$. If some cycle of $N$ meets
no path from $s$ to $t$, then $P_{st}(N)\le2^{b-1}$.
\end{lemma}

\begin{proof}
In a join, $s$ and $t$ have odd degree at most two, and every other vertex has even degree at most
three. So the degrees are one at $s$ and $t$ and zero or two elsewhere. If a cycle $Z$ meets no
path from $s$ to $t$, then $J\mapsto J+Z$ maps the paths injectively to joins that are not paths.
\end{proof}

\begin{lemma}[Reduction]\label{lem:reduction}
Let $N$ be a pole of rank $b$, and let $s,t$ be two of its ports with $P_{st}(N)>2^{b-1}$.
\begin{enumerate}
\item[\textup{(a)}] Replace each port other than $s,t$, with its two edges, by one edge. No loop
arises. The result $N^{\circ}$ is a $2$-pole with ports $s,t$. It has the same rank, the same number
of paths from $s$ to $t$, and the same number of cycles as $N$.
\item[\textup{(b)}] Every bridge of $N^{\circ}$ separates $s$ from $t$ and is incident with neither.
Let $N_1,\dots,N_q$ be the components of $N^{\circ}$ without its bridges, in the order from $s$ to
$t$. Each $N_i$ is a bridgeless $2$-pole of rank $b_i\ge1$, and
\[
  b=\sum_ib_i,\qquad P_{st}(N)=\prod_iP(N_i),\qquad c(N)=\sum_ic(N_i),
\]
where $P(N_i)$ is the number of paths between the two ports of $N_i$. In particular
$P(N_i)\,2^{-b_i}\ge P_{st}(N)\,2^{-b}$ for every $i$.
\end{enumerate}
\end{lemma}

\begin{proof}
(a) Let $v$ be a port other than $s,t$. Suppose that both edges at $v$ end in the same vertex $u$.
Then $u$ has degree three, its third edge is a bridge, and no path from $s$ to $t$ meets the cycle
through $u$ and $v$. This contradicts Lemma~\ref{lem:joins}. So the two neighbours of $v$ are
distinct, and we may replace $v$ by an edge. Paths and cycles correspond. Repeat.

(b) Let $e$ be a bridge of $N^{\circ}$, and let $Y$ be a component of $N^{\circ}-e$, with $n$
vertices, $j$ of them ports. Counting edge ends gives $\beta(Y)=(n-j+1)/2$. This is a positive
integer, so $Y$ contains a cycle. If $Y$ contains neither port or both ports, a cycle on one side
of $e$ meets no path from $s$ to $t$. If $e$ is incident with $s$, every path from $s$ to $t$
starts with $e$ and avoids the cycle on the side of $s$. Both cases contradict
Lemma~\ref{lem:joins}. So every path from $s$ to $t$ passes all bridges, in the same order.

The multigraph $N_i$ is connected and bridgeless. Its vertices have degree three, except the ends
of the two bridges at $N_i$, or $s$, or $t$. These two vertices are distinct: a common end of two
bridges would have a third edge that is a bridge and does not separate $s$ from $t$. So $N_i$ is a
$2$-pole, and it contains a cycle. Ranks add over the $N_i$, every cycle lies in one $N_i$, and a
path from $s$ to $t$ is a concatenation of paths between the ports of the $N_i$ and of the bridges.
The last statement holds because $P(N_j)\le2^{b_j}$ for all $j$.
\end{proof}

\subsection{The shape of a pole with many paths}\label{sec:shape}

Let $N$ be a bridgeless $2$-pole with ports $s,t$ and rank $b\ge2$. It has $2b$ vertices. It has no
cut vertex: a cut vertex of a subcubic multigraph is joined to some component by a single edge.
We use the following notation. The \emph{layer} $L_j$ is the set of vertices at distance $j$ from
$s$. Every edge joins vertices of the same layer or of consecutive layers. Let $\ell$ be the
distance from $s$ to $t$, let $A=w_0w_1\cdots w_\ell$ be a shortest path from $s$ to $t$, and put
$J=N-V(A)$. Then $w_j\in L_j$. A vertex of $A$ lies on two edges of $A$, or on one if it is a
port. So it has at most one \emph{further edge}. This edge leads to $J$, or it is parallel to an
edge of $A$, because $A$ is a shortest path. For $0<j<\ell$ the layer $L_j$ separates $s$ from $t$,
so $|L_j|\ge2$. Call $L_j$ \emph{thin}
if $0<j<\ell$ and $|L_j|=2$. A thin layer consists of $w_j$ and one vertex of $J$. Call a thin
layer \emph{ordinary} if its vertex of $J$ is on no pair of parallel edges
(Figure~\ref{fig:layers}).

\begin{figure}[ht]
\centering
\begin{tikzpicture}[x=1.22cm,y=1cm]
  % layers
  \foreach \j in {1,2,3,6,7,8}
    \draw[slate!45, line width=0.4pt, dash pattern=on 2pt off 2pt] (\j,-0.45) -- (\j,1.75);
  \foreach \j in {1,2,3,6,7,8} \node[note] at (\j,2.0) {$L_{\j}$};
  \node[note] at (0,2.0) {$L_0$};
  \node[note] at (9,2.0) {$L_\ell$};
  % the part of bounded size
  \draw[draw=warm, fill=warm!8, line width=0.6pt, rounded corners=7pt] (3.55,-0.4) rectangle (5.45,1.7);
  \node[note, text=warm, align=center] at (4.5,0.85) {not thin:\\ bounded size};
  % shortest path A
  \foreach \j in {0,...,9} \coordinate (w\j) at (\j,0);
  \foreach \j in {1,2,3,6,7,8} \coordinate (y\j) at (\j,1.3);
  \draw[cool] (w0)--(w1)--(w2)--(w3)--(3.55,0) (5.45,0)--(w6)--(w7)--(w8)--(w9);
  \draw[cool, dash pattern=on 2.5pt off 2pt] (3.55,0)--(5.45,0);
  % edges of J and between
  \draw[edge] (w0)--(y1) (y1)--(y2) (w1)--(y2) (y1)--(w2) (y2)--(y3) (w3)--(y3) (y3)--(3.55,1.3);
  \draw[edge] (5.45,1.3)--(y6) (w6)--(y6) (y6)--(y7) (y7)--(y8) (w7)--(y8) (y7)--(w8) (y8)--(w9);
  \foreach \j in {1,2,3,6,7,8} \node[vtx] at (w\j) {};
  \foreach \j in {1,2,3,6,7,8} \node[vtx] at (y\j) {};
  \node[port] at (w0) {};
  \node[port] at (w9) {};
  \node[note, text=accent, below=3pt] at (w0) {$s$};
  \node[note, text=accent, below=3pt] at (w9) {$t$};
  \node[note, text=accent, below=2pt] at (2,0) {$w_2$};
  \node[note, text=accent, anchor=west] at (9.25,0) {the shortest path $A$};
  \node[note, anchor=west] at (9.25,1.3) {the rest $J$};
\end{tikzpicture}
\caption{A bridgeless $2$-pole with many paths from $s$ to $t$ (schematic). Almost every layer
is thin: it has one vertex of the shortest path $A$ and one vertex of $J$
(Proposition~\ref{prop:thin}).}
\label{fig:layers}
\end{figure}

\begin{idea}{outline of the proof of Proposition~\ref{prop:thin}}
Add the shortest path $A$ to a path from $s$ to $t$. The result is an even subgraph whose edges
away from $A$ form a forest. So more than half of the outside parts are forests, and every layer
is small by Lemma~\ref{lem:width}. By
Lemma~\ref{lem:smallcuts} there are few short cycles away from $A$. By a count of vertices,
almost every layer then has exactly two vertices.
\end{idea}

\begin{proposition}[Thin layers]\label{prop:thin}
For every $\varepsilon\in(0,\half]$ there is a number $C=C_\varepsilon\ge1$ with the following
property. Let $N$ be a bridgeless $2$-pole of rank $b\ge2$ with ports $s,t$ and
$P_{st}(N)\ge(\half+\varepsilon)2^{b}$. Then, in the notation above,
\begin{enumerate}
\item[\textup{(i)}] $\beta(J)\le C$ and $b-C\le\ell\le b$;
\item[\textup{(ii)}] at most $2C$ of the layers $L_j$ with $0<j<\ell$ are not thin, and
$L_j=\emptyset$ for $j>b+C$;
\item[\textup{(iii)}] at most $10C+2$ vertices of $N$ lie outside the ordinary layers.
\end{enumerate}
\end{proposition}

\begin{proof}
\step{Step 1: the layers are small.} Put $P=V(A)$. Then $N[P]$ is connected, and every vertex
outside $P$ has degree three. For a path $Q$ from $s$ to $t$, the set $Q+A$ is an even subgraph of
$N$. Its edges in $N-P$ are edges of the path $Q$, so they form a forest. Different paths give
different even subgraphs, and $N$ has $2^{b}$ even subgraphs. Hence
\begin{equation}\label{eq:qH}
  q(N,P)\;\ge\;\frac{P_{st}(N)}{2^{b}}\;\ge\;\frac12+\varepsilon .
\end{equation}
Let $\lambda(u)$ be the distance from $s$ to $u$ in $N$. This is a level function, and $P$ has at
most one vertex of each level. By Lemma~\ref{lem:width} with $\delta=\varepsilon/2$, every layer
$L_j$ has fewer than $\Theta:=\Theta_{\varepsilon/2}(1)$ vertices of $J$.

\step{Step 2: a short cycle.} Put $U_j=L_j\sm V(A)$, and for an interval $I$ of indices let $J_I$ be
the subgraph of $J$ induced by the sets $U_j$ with $j\in I$. Suppose that $J_I$ contains a cycle.
Choose a cycle $Z$ of $J_I$ such that the interval $[a,a']$ of the layers that $Z$ meets is as
short as possible, and put $\sigma=a'-a$ and $R=U_a\cup\dots\cup U_{a'}$. We claim that
$|R|\le20\Theta$.

By the choice of $Z$, every cycle of $J[R]$ meets $U_a$. So $J[R]-U_a$ is a forest. Deleting a
vertex of degree at most three lowers the rank by at most two, so $\beta(J[R])\le2|U_a|<2\Theta$.
The vertices of $R$ have degree three in $N$. Counting edge ends gives
$|R|=2\beta(J[R])+|\delta R|-2\comp(J[R])$. An edge that leaves $R$ ends in $U_{a-1}$, in
$U_{a'+1}$, or in one of the vertices $w_{a-1},\dots,w_{a'+1}$, each of which has at most one
further edge. So $|\delta R|<6\Theta+\sigma+3$, and $|R|<10\Theta+\sigma+1$. The cycle $Z$ has a vertex in
layer $a$ and a vertex in layer $a'$, so $|Z|\ge2\sigma$. Also $|Z|\le|R|$. Hence
$\sigma\le10\Theta$ and $|R|\le20\Theta$.

\step{Step 3: the rank of $J$.} Start with the interval of all layers. As long as $J_I$ contains a
cycle for some interval $I$ of the present family, choose $Z$ and $R$ as in Step~2, and replace $I$
by the at most two intervals that remain when $[a,a']$ is removed. Let $Z_1,\dots,Z_k$ be the
chosen cycles and $R_1,\dots,R_k$ the corresponding sets. The sets $R_i$ are pairwise disjoint,
$|R_i|\le20\Theta$, and $J-\bigcup_iR_i$ is a forest, because an edge of $J$ joins equal or
consecutive layers. So $\beta(J)\le2\sum_i|R_i|\le40\Theta k$.

The sets $V(Z_i)$ are pairwise disjoint cyclic regions of $N$ that avoid $P$. Each has at most
$|Z_i|\le20\Theta$ edges that leave it. If $\xi\cap E(N-P)$ is a forest, then
$\xi\cap E(N[V(Z_i)])$ is a forest for every $i$. So Lemma~\ref{lem:smallcuts} and~\eqref{eq:qH}
give $k\le(20\Theta+1)2^{20\Theta-1}/\varepsilon$. Hence $\beta(J)\le C$ with
\[
  C=C_\varepsilon:=\frac{40\Theta\,(20\Theta+1)\,2^{20\Theta-1}}{\varepsilon}.
\]

\step{Step 4: counting.} The layers $L_0,\dots,L_\ell$ contain at least $2\ell$ vertices, so
$\ell\le b$. The vertices of $A$ meet the $\ell$ edges of $A$ and at most $\ell+1$ other edges.
If $J$ is not empty, then $|E(J)|\ge(3b-1)-(2\ell+1)$ and $|V(J)|=2b-\ell-1$, so
$\beta(J)\ge|E(J)|-|V(J)|+1\ge b-\ell$. If $J$ is empty, then $2b=\ell+1$. In both cases
$\Delta:=b-\ell\le C$. This is~(i).

Compare $N$ with a multigraph that has one vertex in $L_0$ and in $L_\ell$, two in each layer
between, and none beyond. Then $N$ has exactly $2\Delta$ vertices more. So at most $2\Delta$
layers between $L_0$ and $L_\ell$ are not thin, and at most $2\Delta$ layers lie beyond $L_\ell$.
This is~(ii). A layer between $L_0$ and $L_\ell$ with $w\ge3$ vertices has $w-2\ge w/3$ vertices
too many. So at most $2+3\cdot2\Delta$ vertices lie outside the thin layers.

Finally let $y\in J$ be on a pair of parallel edges. The other end of the pair has two edges to
$y$, so it is not on $A$, whose vertices have at most one further edge. So the pair is a cycle of
$J$. Two such pairs have no common vertex,
because $N$ is subcubic and has more than two vertices. So there are at most $\beta(J)\le C$ pairs,
and at most $2C$ thin layers are not ordinary. They contain at most $4C$ vertices. This gives~(iii).
\end{proof}

\begin{theorem}[Few cycles]\label{thm:fewcycles}
For every $\varepsilon\in(0,\half]$ there is a constant $K_\varepsilon$ with the following property.
Let $N$ be a pole of rank $b$, and let $s,t$ be two of its ports with
$P_{st}(N)\ge(\half+\varepsilon)2^{b}$. Then
\[
  c(N)\;\le\;K_\varepsilon\,2^{b/2}.
\]
\end{theorem}

\begin{proof}
Let $C=C_\varepsilon$ be as in Proposition~\ref{prop:thin}. We prove that a bridgeless $2$-pole
with these properties has at most $2^{31C+22}\,2^{b/2}$ cycles. The theorem follows with
$K_\varepsilon=2^{31C+23}$. Indeed, in the notation of Lemma~\ref{lem:reduction}, each $N_i$ has the
required number of paths, and $\sum_i2^{b_i/2}\le2^{b/2}+q-1\le2\cdot2^{b/2}$, because
$x+y\le xy+1$ for $x,y\ge1$.

So let $N$ be a bridgeless $2$-pole. If $b=1$, then $N$ has one cycle. Let $b\ge2$, and use the
notation above. Let $Z$ be a cycle, and let $L_i$ and $L_j$ be the first and the last layer that it
meets. Let $M$ be the set of the vertices that lie in $L_i$, in $L_j$, or outside the ordinary
layers. By Proposition~\ref{prop:thin}(iii), $|M|\le10C+6$.

\step{Two strands.} Let $L_l=\{w_l,y_l\}$ be an ordinary layer with $i<l<j$. The cycle $Z$ has
vertices before and after this layer. So it has at least two edges from $L_{l-1}$ to $L_l$ and at
least two edges from $L_l$ to $L_{l+1}$. The two vertices of $L_l$ have degree at most two in $Z$.
So $Z$ has exactly two edges on each side of $L_l$ and does not contain an edge $w_ly_l$. Suppose
that both edges from $L_{l-1}$ end in the same vertex of $L_l$. Then both edges to $L_{l+1}$ start
in the other vertex, and no edge of $Z$ joins the part of $Z$ before $L_l$ to the part after it.
This is impossible, because $Z$ is connected. So each vertex of $L_l$ has one edge of $Z$ to
$L_{l-1}$ and one to $L_{l+1}$.

\step{Matchings.} Let $L_l$ and $L_{l+1}$ be ordinary layers with $i<l$ and $l+1<j$. By the last
paragraph, the edges of $Z$ between them form a perfect matching between $\{w_l,y_l\}$ and
$\{w_{l+1},y_{l+1}\}$. We count the perfect matchings. The edge $y_ly_{l+1}$ is not doubled,
because the layers are ordinary. The edges $w_ly_{l+1}$ and $y_lw_{l+1}$ are not doubled, because
a vertex of $A$ has at most one further edge. If the edge $w_lw_{l+1}$ of $A$ is doubled, then
$w_l$ and $w_{l+1}$ have no edge to $J$. So there are at most two perfect matchings. If there
are two, then $w_{l+1}$ has two edge ends towards $L_l$: either both $w_l$ and $y_l$ are its
neighbours, or $w_lw_{l+1}$ is doubled. In the same way $w_{l+1}$ has two edge ends towards
$L_{l+2}$ if there are two perfect matchings between $L_{l+1}$ and $L_{l+2}$. Both cannot happen,
because $w_{l+1}$ has degree three. So among consecutive pairs of ordinary layers, no two
neighbouring pairs have two perfect matchings each.

\step{The count.} An edge of $Z$ has an end in $M$, or it joins two consecutive ordinary layers
strictly between $L_i$ and $L_j$. So $Z$ is determined by $i$ and $j$, by its edges at $M$, and by
its perfect matchings. There are at most $2^{3|M|}$ sets of edges at $M$. By the last paragraph
there are at most $2^{(j-i)/2}$ choices for the matchings. All layers have an index of at most
$n:=b+C$, by Proposition~\ref{prop:thin}(ii). Hence
\[
  c(N)\;\le\;2^{30C+18}\sum_{0\le i\le j\le n}2^{(j-i)/2}
  \;\le\;2^{30C+18}\cdot12\cdot2^{(b+C)/2}\;\le\;2^{31C+22}\,2^{b/2}.\qedhere
\]
\end{proof}

A pole with many paths also has a short cycle:

\begin{lemma}[Four thin layers]\label{lem:fourlayers}
Let $N$ be a bridgeless $2$-pole of rank at least two. If four consecutive layers are thin, then
$N$ has a cycle of length at most four all of whose vertices lie in these four layers.
\end{lemma}

\begin{proof}
The eight vertices of the four layers have degree three, and each has a neighbour in the layer
before its own. The edges between two consecutive layers form a cut of $N$, so there are at least
two of them. There are at most four: the two vertices of the earlier layer have six edge ends, and
two of these point backwards. A pair of parallel edges among the eight vertices is a cycle of
length two. Suppose there is none. If some cut between two of the four layers has four edges, it
is a cycle of length four. So let the three cuts have $e_1,e_2,e_3\in\{2,3\}$ edges. Let $r_2$ and
$r_3$ be the numbers of edges inside the second and the third layer. Counting edge ends in these
layers gives $e_1+e_2+2r_2=6$ and $e_2+e_3+2r_3=6$. So $e_1=e_2=e_3$.

If all three are $2$, then $r_2=r_3=1$, and each vertex of the second and third layer has one edge
to each side. The two edges between the second and the third layer and the edges inside these
layers form a cycle of length four. If all three are $3$, then $r_2=0$. In the second layer one
vertex $u$ has two neighbours in the first layer and the other vertex $u'$ has one, say $x$; then
$x$ is also a neighbour of $u$. So $u$ has one neighbour $y$ in the third layer and $u'$ has two,
among them $y$. Then $xuyu'$ is a cycle of length four.
\end{proof}

\begin{corollary}\label{cor:nodense}
For every $\varepsilon\in(0,\half]$ and every integer $g\ge0$ there is $b_1(\varepsilon,g)$ with the
following property. Let $G$ be a connected cubic multigraph without loops, of rank at least $9$,
that has no cyclic cut of size at most four. Let $K$ be an induced subgraph of $G$ that is a pole
with at most $g+2$ ports and rank $\beta(K)\ge b_1(\varepsilon,g)$. Then
$P_{st}(K)<(\half+\varepsilon)2^{\beta(K)}$ for all ports $s,t$ of $K$.
\end{corollary}

\begin{proof}
Suppose that $P_{st}(K)\ge(\half+\varepsilon)2^{\beta(K)}$, and let $N=K^{\circ}$ and
$N_1,\dots,N_q$ be as in Lemma~\ref{lem:reduction}. An edge of $N$ is a path of $K$ whose inner
vertices are ports of $K$. Call the edge \emph{plain} if this path has no inner vertex. At most
$g$ edges of $N$ are not plain.

Let $1<i<q$, and suppose that all edges of $N_i$ are plain. Then the vertices of $N_i$ have degree
three in $K$. So exactly two edges of $G$ leave $V(N_i)$, one for each bridge at $N_i$. The
multigraph $N_i$ contains a cycle, and so does $N_1$. This is a cyclic cut of size two, which $G$
does not have. So $q\le g+2$, and some $N_i$ has rank $b'\ge\beta(K)/(g+2)$.

Let $C=C_\varepsilon$. By Proposition~\ref{prop:thin}, applied to $N_i$, the distance between its
ports is at least $b'-C$, and at most $2C$ of its inner layers are not thin. At most $2g$ layers
contain an end of an edge that is not plain. Divide the inner layers into groups of four
consecutive layers. If $b'\ge9C+8g+9$, then in some group all layers are thin and all edges with
both ends in the group are plain. Lemma~\ref{lem:fourlayers} gives a cycle $Z$ of length at most
four in this group. It is a cycle of $G$. At most four edges of $G$ leave $V(Z)$, and at most
eight edges of $G$ have an end in $V(Z)$. So $G-V(Z)$ has rank at least $\beta(G)-8>0$, and
$\delta V(Z)$ is a cyclic cut of size at most four. This is a contradiction. So the corollary
holds with $b_1(\varepsilon,g)=(g+2)(9C+8g+9)$.
\end{proof}

\paragraph{Consequences of Theorem~\ref{thm:fewcycles}.}
From now on fix
\[
  \varepsilon_0=\tfrac1{1000},\qquad K_0=K_{\varepsilon_0},\qquad\kappa(h)=2K_0\,2^{-h/2}.
\]
For four distinct ports $a,b,c,d$ of a pole $N$ let $f_{ab|cd}(N)$ be the number of pairs of
vertex-disjoint paths of $N$, one from $a$ to $b$ and one from $c$ to $d$.

\begin{lemma}\label{lem:gadgets}
Let $N$ be a pole of rank $h$. In each statement below, the letters denote distinct ports of $N$.
Write $P_{xy}$ for $P_{xy}(N)$ and $f$ for $f(N)$.
\begin{enumerate}
\item[\textup{(a)}] If $P_{ab}\ge0.501\cdot2^{h}$, then $c(N)\le K_0\,2^{h/2}$.
\item[\textup{(b)}] If $P_{ac}+P_{bc}\ge1.002\cdot2^{h}$, then $P_{ab}\le\kappa(h)\,2^{h}$.
\item[\textup{(c)}] If $P_{ac}+P_{ad}+P_{bc}+P_{bd}\ge2.004\cdot2^{h}$, then
$f_{ac|bd}+f_{ad|bc}\le\kappa(h)\,2^{h}$.
\item[\textup{(d)}] If $P_{ab}+f_{ac|bd}+f_{ad|bc}\ge1.002\cdot2^{h}$, then
$P_{cd}\le\kappa(h)\,2^{h}$.
\end{enumerate}
\end{lemma}

\begin{proof}
(a) is Theorem~\ref{thm:fewcycles}. In each of the other cases we change $N$ into a pole $N'$ and
apply Theorem~\ref{thm:fewcycles} to two ports of $N'$.

(b) Add a new vertex $z$ and join it to $a$ and $b$. Then $N'$ has rank $h+1$, and
$P_{zc}(N')=P_{ac}+P_{bc}\ge(\half+\varepsilon_0)2^{h+1}$. The cycles of $N'$ through $z$
correspond to the paths of $N$ from $a$ to $b$. So $P_{ab}\le c(N')\le K_0\,2^{(h+1)/2}$.

(c) Add a new vertex $z$ joined to $a$ and $b$, and a new vertex $w$ joined to $c$ and $d$. Then
$N'$ has rank $h+2$, and $P_{zw}(N')$ is the sum of the four path counts. The cycles of $N'$
through $z$ and $w$ correspond to the pairs counted by $f_{ac|bd}+f_{ad|bc}$. So this sum is at
most $K_0\,2^{(h+2)/2}$.

(d) Add an edge between $c$ and $d$. Then $N'$ has rank $h+1$. A path from $a$ to $b$ in $N'$
avoids the new edge, or it consists of the new edge and a pair counted by
$f_{ac|bd}+f_{ad|bc}$. So $P_{ab}(N')=P_{ab}+f_{ac|bd}+f_{ad|bc}$. The cycles of $N'$ through the
new edge correspond to the paths of $N$ from $c$ to $d$. So $P_{cd}\le K_0\,2^{(h+1)/2}$.
\end{proof}

\subsection{Full pairs}\label{sec:full}

A \emph{square} is a copy of the complete bipartite graph $K_{2,2}$; its two classes are its
\emph{left pair} and its \emph{right pair}. A \emph{single rung} is one edge; its two ends form
both its left pair and its right pair. These are the two kinds of \emph{tiles}. The \emph{weight}
of a square is $2$, and the weight of a single rung is $1$. A \emph{chain} of tiles
$\tau_1,\dots,\tau_m$ consists of these tiles, pairwise disjoint, and of two disjoint edges from
the right pair of $\tau_i$ to the left pair of $\tau_{i+1}$ for $1\le i<m$. These two edges are an
\emph{interface}. A \emph{capped ladder} is a chain with two more vertices $s$ and $t$, the
\emph{caps}: $s$ is joined to the two vertices of the left pair of $\tau_1$, and $t$ to the two
vertices of the right pair of $\tau_m$ (Figure~\ref{fig:ladder}). A capped ladder is a $2$-pole. If
its tiles have total weight $W$, it has $2W+2$ vertices and rank $W+1$.

\begin{figure}[ht]
\centering
\begin{tikzpicture}[x=1cm,y=1cm]
  \coordinate (s) at (0,0);
  \coordinate (t) at (9.6,0);
  \foreach \n/\x in {a/1.2,b/2.4,c/3.8,d/5.2,e/6.4,f/7.8}{
    \coordinate (\n u) at (\x,0.65);
    \coordinate (\n l) at (\x,-0.65);}
  \draw[edge] (s)--(au) (s)--(al);
  \draw[cool] (au)--(bu) (al)--(bl) (au)--(bl) (al)--(bu);
  \draw[edge] (bu)--(cu) (bl)--(cl);
  \draw[hot] (cu)--(cl);
  \draw[edge] (cu)--(du) (cl)--(dl);
  \draw[cool] (du)--(eu) (dl)--(el) (du)--(el) (dl)--(eu);
  \draw[edge] (eu)--(fu) (el)--(fl);
  \draw[hot] (fu)--(fl);
  \draw[edge] (fu)--(t) (fl)--(t);
  \foreach \q in {au,al,bu,bl,cu,cl,du,dl,eu,el,fu,fl} \node[vtx] at (\q) {};
  \node[port] at (s) {};
  \node[port] at (t) {};
  \node[note, text=accent, left=3pt] at (s) {$s$};
  \node[note, text=accent, right=3pt] at (t) {$t$};
  \node[note, text=accent] at (1.8,-1.15) {square};
  \node[note, text=warm] at (3.8,-1.15) {single rung};
  \node[note, text=accent] at (5.8,-1.15) {square};
  \node[note, text=warm] at (7.8,-1.15) {single rung};
  \node[note] at (3.1,1.1) {interface};
  \draw[slate, line width=0.4pt] (3.1,0.95) -- (3.1,0.7);
\end{tikzpicture}
\caption{A capped ladder with four tiles, of weight $6$ and rank $7$. Without $s$ and $t$ it is a
chain. If the right pair of the last tile is joined to the left pair of the first tile instead, the
result is a ring (Section~\ref{sec:rings}).}
\label{fig:ladder}
\end{figure}

\begin{lemma}[Ladders]\label{lem:ladders}
Let $N$ be a bridgeless $2$-pole whose pair of ports is full. Then $N$ consists of two vertices and
two parallel edges, or $N$ is a capped ladder whose caps are the ports.
\end{lemma}

\begin{proof}
Let $s,t$ be the ports. By Lemma~\ref{lem:joins}, every cycle meets every path from $s$ to $t$. If
the rank is one, $N$ is a cycle with two vertices. Otherwise $N$ has no cut vertex, so there are
two paths $Q$ and $Q'$ from $s$ to $t$ without a common inner vertex. Let $F$ be the set of the
edges outside $Q\cup Q'$. In $F$ the inner vertices of $Q$ and $Q'$ have degree one, the other
vertices except $s,t$ have degree three, and $s,t$ have degree zero.

A cycle in $F$ would avoid $Q$. So $F$ is a forest. Suppose that a component of $F$ has two leaves
on $Q$. The path of $F$ between them and a part of $Q$ form a cycle that avoids $Q'$. This is
impossible, and the same holds for $Q'$. So every component of $F$ has at most two leaves. A tree
with degrees one and three has two more leaves than vertices of degree three. So every component
of $F$ with an edge is a single edge, from an inner vertex of $Q$ to an inner vertex of $Q'$. Hence
all vertices lie on $Q\cup Q'$, and $F$ is a perfect matching between the inner vertices
$x_1,\dots,x_k$ of $Q$ and $y_1,\dots,y_k$ of $Q'$, in the order from $s$ to $t$. Write its edges
as $x_iy_{\pi(i)}$.

Let $i<j<l$. If $\pi(l)<\min\{\pi(i),\pi(j)\}$, then the matching edges at $x_i$ and $x_j$ and the
parts of $Q$ and $Q'$ between their ends form a cycle. The path that follows $Q'$ from $s$ to
$y_{\pi(l)}$, crosses to $x_l$ and follows $Q$ to $t$ avoids it. If
$\pi(i)>\max\{\pi(j),\pi(l)\}$, the same holds for the cycle through the matching edges at $x_j$
and $x_l$ and the path that follows $Q$ to $x_i$, crosses, and follows $Q'$ to $t$. So neither
case occurs. It follows that $\pi(1)\le2$, and that $\pi(2)=1$ if $\pi(1)=2$. By induction, $\pi$
is composed of fixed points and of exchanges of two consecutive numbers. A fixed point $i$ is a
single rung with the ends $x_i,y_i$. An exchange of $i$ and $i+1$ is a square with the left pair
$x_i,y_i$ and the right pair $x_{i+1},y_{i+1}$. The remaining edges of $Q$ and $Q'$ are the
interfaces and the edges at the caps.
\end{proof}

We need two counts in a capped ladder. Let $\tau_1,\dots,\tau_m$ be its tiles. First, let
$1\le i\le m$ and let $u$ be a vertex of the left pair of $\tau_i$. The number of paths from $u$ to
$t$ inside $\tau_i\cup\dots\cup\tau_m\cup\{t\}$ is $2^{q}$, where $q$ is the weight of
$\tau_i,\dots,\tau_m$. Indeed, such a path passes the tiles in order. It enters a tile at a vertex
of the left pair and leaves it at a vertex of the right pair. In a square there are two paths to
each vertex of the right pair. In a single rung the path stays or crosses. Second, let $A_i$ be
the number of paths between the two vertices of the right pair of $\tau_i$ inside
$\{s\}\cup\tau_1\cup\dots\cup\tau_i$, and $A_0=1$. This multigraph has rank equal to the weight
$p$ of $\tau_1,\dots,\tau_i$. So $A_i\le2^{p}$ by~\eqref{eq:joinbound}.

\begin{lemma}[A third port]\label{lem:thirdport}
Let $T$ be a $3$-pole of rank $h$ with ports $s,t,v$. If the pair $s,t$ is full, then
$P_{vt}(T)=2^{h}$ or $P_{vt}(T)\le\frac78\,2^{h}$.
\end{lemma}

\begin{proof}
By Lemma~\ref{lem:reduction}, $T$ arises from a $2$-pole $N$ with ports $s,t$ when an edge $e$ is
subdivided by $v$, and $N$ consists of bridges and of bridgeless $2$-poles $N_1,\dots,N_q$ with
full pairs. If $e$ is a bridge or an edge of some $N_i$ with $i\ge2$, then no path from $v$ to $t$
meets $N_1$. Such a path lies in a part of $T$ of rank at most $h-1$, so $P_{vt}(T)\le2^{h-1}$
by~\eqref{eq:joinbound}. Let $e$ be an edge of $N_1$, and let $t_1$ be the second port of $N_1$. A
path from $v$ to $t$ is a path from $v$ to $t_1$ followed by paths through $N_2,\dots,N_q$, and
each of these $2$-poles is full. So we may assume that $q=1$.

By Lemma~\ref{lem:ladders}, $N$ is a cycle with two vertices or a capped ladder. In the first
case $T$ is a triangle and $P_{vt}=2=2^{h}$. In the second case let $p$ be the weight of the
tiles before $e$ and $q'$ the weight of the tiles after $e$; if $e$ lies in a tile, this tile is
not counted. A path from $v$ to $t$ never returns once it has entered the tiles after $e$, and
from its point of entry it has $2^{q'}$ continuations. There are three cases.

\step{The edge $e$ is in an interface or at a cap.} Then $h=p+q'+1$. The path goes from $v$ to
the right, or it goes to the left, follows one of the $A$ paths between the two ends of the
interface on the left side, and uses the other edge of the interface. Here $A=A_i$ for the tile
$\tau_i$ before $e$. So $P_{vt}=(1+A)2^{q'}$. If $p=0$, this is $2^{h}$. If $p\ge1$, then
$A\le2^{p}$ and $P_{vt}\le\frac34\,2^{h}$.

\step{The edge $e$ is a single rung.} Then $h=p+q'+2$. From each end of the rung the path
continues to the right, or it makes the detour through the tiles before the rung to the other end
of the rung and continues from there. So $P_{vt}=2(1+A)2^{q'}$, and the conclusion is the same.

\step{The edge $e$ is in a square.} Then $h=p+q'+3$. Let $u,u'$ be the left pair and $x,x'$ the
right pair of the square, and let $e=ux$. A path from $v$ leaves the square at $x$ or $x'$.
Without a detour through the tiles before the square, its part in the square is one of
\[
  vx,\qquad vux'u'x,\qquad vux',\qquad vxu'x' .
\]
With a detour between $u$ and $u'$ it is one of $vu\cdots u'x$, $vu\cdots u'x'$ and
$vxu'\cdots ux'$. There are $A$ detours. So $P_{vt}=(4+3A)2^{q'}$. If $p=0$, then $A=1$ and
$P_{vt}=\frac78\,2^{h}$. If $p\ge1$, then $P_{vt}\le(4+3\cdot2^{p})2^{q'}\le\frac58\,2^{h}$.
\end{proof}

\begin{theorem}[Gap]\label{thm:gap}
Let $N$ be a pole of rank $b$, and let $s,t$ be two of its ports. Then
\[
  P_{st}(N)=2^{b}\qquad\text{or}\qquad P_{st}(N)\le\tfrac{15}{16}\,2^{b}.
\]
\end{theorem}

\begin{proof}
We use induction on $b$. Let $P_{st}(N)>2^{b-1}$. By Lemma~\ref{lem:reduction} we may assume that
$N$ is a $2$-pole. If $N$ has a bridge, the claim holds for the multigraphs $N_i$ of
Lemma~\ref{lem:reduction}, which have smaller rank, and the path counts multiply. So let $N$ be
bridgeless. If $b=1$, the pair is full. Let $b\ge2$. If $s$ and $t$ are adjacent, then
$P_{st}(N)\le1+2^{b-1}\le\frac34\,2^{b}$, by~\eqref{eq:joinbound} for $N$ without that edge.
Otherwise let $x$ and $y$ be the neighbours of $s$. They are distinct, because $N$ has no bridge.
The multigraph $T=N-s$ is connected, because $N$ has no cut vertex. It is a $3$-pole of rank $b-1$
with ports $x,y,t$, and
\[
  P_{st}(N)=P_{xt}(T)+P_{yt}(T).
\]
By induction each term is $2^{b-1}$ or at most $\frac{15}{16}2^{b-1}$. If neither pair is full, we
are done. If the pair $x,t$ is full, then $P_{yt}(T)$ is $2^{b-1}$ or at most $\frac78\,2^{b-1}$,
by Lemma~\ref{lem:thirdport}. So $P_{st}(N)$ is $2^{b}$ or at most $\frac{15}{16}2^{b}$.
\end{proof}

\begin{corollary}[Four ports]\label{cor:fourports}
Let $N$ be a $4$-pole with ports $a,b,c,d$ such that the four pairs $ac$, $ad$, $bc$, $bd$ are
full. Then one of the following holds.
\begin{enumerate}
\item[\textup{(i)}] $N$ is a chain of tiles; $a,b$ is the left pair of its first tile and $c,d$ is
the right pair of its last tile.
\item[\textup{(ii)}] $N$ has a bridge $e$ such that one component of $N-e$ contains $a$, $b$ and a
cycle, and the other component contains $c$, $d$ and a cycle.
\end{enumerate}
\end{corollary}

\begin{proof}
Add a vertex $z$ joined to $a,b$ and a vertex $w$ joined to $c,d$. The result $N'$ is a $2$-pole
of rank $\beta(N)+2$ with ports $z,w$, and $P_{zw}(N')=4\cdot2^{\beta(N)}$. So the pair $z,w$ is
full. Let $N_1,\dots,N_q$ be as in Lemma~\ref{lem:reduction}. Each $N_i$ has a full pair of ports,
so Lemma~\ref{lem:ladders} applies to it. The neighbours $a,b$ of $z$ are distinct, so $N_1$ is a
capped ladder with the cap $z$, and $a,b$ is the left pair of its first tile. In the same way
$N_q$ is a capped ladder with the cap $w$. If $q=1$, this is~(i). If $q\ge2$, let $e$ be the
bridge at $N_1$. The multigraphs $N_1-z$ and $N_q-w$ are connected, each contains a cycle, and they
lie in different components of $N-e$. This is~(ii).
\end{proof}

\section{Rings, small cyclic cuts, and cycles of length four}\label{sec:cuts}

In this section $G$ is a connected cubic multigraph of rank $r$.

\begin{idea}{the results of this section}
We count the cycles of the rings. Then we show that a graph of large rank with a cyclic cut of
size at most three has at most $2^{r-1}+O(2^{r/2})$ cycles (Proposition~\ref{prop:threecut}), and
that a graph without such a cut and with a cycle of length four is a ring, or it has fewer than
$2^{r-1}$ cycles (Theorem~\ref{thm:fourcut}). For the second statement we delete the cycle of
length four and close the rest by a \emph{cap}, a small tree
(Lemma~\ref{lem:cap}). The proof uses the gap of Section~\ref{sec:poles}: a pole has all
possible paths between two ends, or a fixed proportion of them is missing.
\end{idea}

Let $X$ be a set of $n\ge1$ vertices of $G$ with $|\delta X|=d$. Counting edge ends gives
$\beta(G[X])=(n-d)/2+\comp(G[X])$. So $G[X]$ contains a cycle if $d\le2$, and also if $d=3$ and
$n\ge2$. Hence every cut of size at most two is cyclic, and a cut of size three is cyclic unless
one of its sides is a single vertex. If $G$ has a loop at a vertex $v$ and $r\ge2$, the other
edge at $v$ is a cyclic cut of size one.

\subsection{Rings}\label{sec:rings}

A \emph{ring} consists of pairwise disjoint tiles $\tau_1,\dots,\tau_m$ and of two disjoint edges
from the right pair of $\tau_i$ to the left pair of $\tau_{i+1}$ for every $i$, with indices modulo
$m$. These pairs of edges are the \emph{interfaces}. Let the ring have $a$ squares and $s$ single
rungs, and let $W=2a+s$ be its weight. We assume that $W\ge7$. Then the ring is a cubic graph with
$2W$ vertices and rank $W+1$. Put
\[
  E(a,s)=2^{a}\Bigl(2a+4s+\tbinom s2+\tfrac12\Bigr)-3a-4s+\tbinom s2 .
\]
Then $2^{2q}+E(q,0)$ and $2^{2q+1}+E(q,1)$ are the numbers $H(2q+1)$ and $H(2q+2)$
of~\eqref{eq:H}.

\begin{proposition}[Rings]\label{prop:rings}
A ring with $a\ge1$ squares and $s$ single rungs, of weight $W\ge7$, has at most $2^{W}+E(a,s)$
cycles. Equality holds if $s\le1$. A ring of weight
$W\ge7$ without squares has at most $2^{W}+W^{2}$ cycles.
\end{proposition}

\begin{figure}[ht]
\centering
\begin{tikzpicture}[x=1cm,y=1cm]
  \foreach \n/\x in {a/1.0,b/2.2,c/3.4,d/4.6,e/5.8,f/7.0,g/8.2}{
    \coordinate (\n u) at (\x,0.6);
    \coordinate (\n l) at (\x,-0.6);}
  \draw[edge, dash pattern=on 2pt off 1.6pt] (0.1,0.6)--(au) (0.1,-0.6)--(al) (gu)--(9.1,0.6) (gl)--(9.1,-0.6);
  \draw[edge] (bu)--(cl) (bl)--(cu);
  \draw[edge] (du)--(eu) (dl)--(el);
  \draw[edge] (eu)--(fu) (el)--(fl) (eu)--(fl) (el)--(fu);
  \draw[edge] (fu)--(gu) (fl)--(gl) (gu)--(gl);
  \draw[hot] (au)--(al) (au)--(bu)--(cu)--(du)--(dl)--(cl)--(bl)--(al);
  \foreach \q in {au,al,bu,bl,cu,cl,du,dl,eu,el,fu,fl,gu,gl} \node[vtx] at (\q) {};
  \node[note, text=warm] at (1.0,-1.05) {$\tau_i$};
  \node[note, text=warm] at (4.6,-1.05) {$\tau_j$};
  \node[note] at (2.8,-1.05) {a square};
  \node[note] at (6.4,-1.05) {a square};
  \node[note] at (8.2,-1.05) {a single rung};
  \node[note, text=warm, anchor=west] at (9.4,0) {the cycle turns in $\tau_i$ and in $\tau_j$};
\end{tikzpicture}
\caption{A part of a ring, and a cycle that uses both edges of every interface from $\tau_i$ to
$\tau_j$ and no other interface. It turns in the single rungs $\tau_i$ and $\tau_j$, and it
passes the square between them on two disjoint edges.}
\label{fig:turn}
\end{figure}

\begin{proof}
Two interfaces together form a cut. So a cycle uses an odd number of edges of every interface, or
an even number of edges of every interface.

\step{One edge of every interface.} There are $2^{m}$ ways to choose the edges. The cycle then
enters each tile at one vertex of its left pair and leaves it at one vertex of its right pair. In a
square there are two such paths. In a single rung there is one: the cycle stays at the vertex or
crosses the rung. Every choice gives a cycle. So there are $2^{m+a}=2^{W}$ cycles of this kind.

\step{No interface.} These cycles are the $a$ squares.

\step{Both edges of every interface.} In a square the cycle then uses two disjoint edges, and there
are two ways to choose them. In a single rung it does not use the rung. Each of the $2^{a}$ choices
is a cycle or a union of two cycles. Changing the choice in one square changes one into the other.
So there are $2^{a-1}$ such cycles if $a\ge1$, and at most one if $a=0$.

\step{Both edges of some interfaces, but not of all.} Let the cycle use the interface after
$\tau_i$ and not the interface before $\tau_i$ (Figure~\ref{fig:turn}). Then it turns in $\tau_i$: it joins the two
vertices of the right pair inside the tile. It passes the next tiles on two disjoint paths, until
it turns in a tile $\tau_j$. This is already a closed curve, so the cycle uses exactly the
interfaces from $\tau_i$ forward to $\tau_j$. In each tile from $\tau_i$ to $\tau_j$ it has two
possibilities if the tile is a square, and one if it is a single rung, in either role. Every
choice gives a cycle. Let $I(i,j)$ be the set of the other tiles, those strictly after $\tau_j$
and strictly before $\tau_i$. Then the number of these cycles is $2^{a}\,\omega(i,j)$, where
$\omega(i,j)=2^{-n}$ and $n$ is the number of squares in $I(i,j)$. Every cycle of the present kind
arises for exactly one ordered pair $(i,j)$ of distinct tiles.

So a ring with $a\ge1$ has $2^{W}+a+2^{a-1}+2^{a}\,\Omega$ cycles, where $\Omega$ is the sum of
$\omega(i,j)$ over the ordered pairs of distinct tiles. We compute $\Omega$. Let $\tau_j$ be a
square. The squares $\tau_i\ne\tau_j$ give $\omega=1,\frac12,\dots,2^{2-a}$, in the order in which
they follow $\tau_j$. So the pairs of two squares contribute $a(2-2^{2-a})$. Let $\tau_i$ be a
single rung. The squares $\tau_j$ give $\omega=1,\frac12,\dots,2^{1-a}$, in the order in which they
precede $\tau_i$. The same holds with the roles of $i$ and $j$ exchanged. So the pairs of a single
rung and a square contribute $2s(2-2^{1-a})$. Finally let $\tau_i,\tau_j$ be two single rungs, with
$n$ squares on one side between them and $a-n$ on the other. The two ordered pairs contribute
$2^{-n}+2^{-(a-n)}\le1+2^{-a}$, with equality if and only if $n=0$ or $n=a$. Hence
\[
  \Omega\;\le\;2a+4s+\tbinom s2-2^{-a}\Bigl(4a+4s-\tbinom s2\Bigr),
\]
with equality if $s\le1$, because then there is no pair of two single rungs. This is the first
statement.

If $a=0$, every ordered pair of distinct tiles gives one cycle. There are $W(W-1)$ pairs, and at
most one further cycle uses both edges of all interfaces.
\end{proof}

\begin{corollary}[Rings of a given rank]\label{cor:ringmax}
Let $r\ge8$. Every ring of rank $r$ has at most $H(r)$ cycles. A ring of rank $r$ with
$\lfloor(r-1)/2\rfloor$ squares and at most one single rung has exactly $H(r)$ cycles.
\end{corollary}

\begin{proof}
Let $W=r-1$. A ring of weight $W$ with $a=\lfloor W/2\rfloor$ squares has $s=W-2a\le1$ single
rungs, and it has $2^{W}+E(a,s)=H(r)$ cycles by Proposition~\ref{prop:rings}. For $a\ge2$ and
$2a+s=W$, a direct computation gives
\[
  E(a-1,s+2)-E(a,s)=-2^{a-1}\Bigl(W+\frac{s^{2}+s}2-\frac{13}2\Bigr)+2s-4\;<\;0,
\]
because $W\ge7$. So $2^{W}+E(a,s)\le H(r)$ for every ring of weight $W$ with a square. The number
$H(r)-2^{W}$ is $71$ for $W=7$ and $124$ for $W=8$, and it more than doubles when $W$ grows by
two. So it exceeds $W^{2}$, and a ring without squares has fewer than $H(r)$ cycles.
\end{proof}

We write $R_r$ for a ring of rank $r\ge8$ with $\lfloor(r-1)/2\rfloor$ squares and at most one
single rung. In the notation of Horak and Kim~\cite{HK} it is the graph $MC_{r-1}$, and they
computed $c(R_r)=H(r)$.

\subsection{Cuts of size at most three}\label{sec:threecut}

In this subsection we use Theorem~\ref{thm:fewcycles} and Lemma~\ref{lem:gadgets} with one
number $\varepsilon_0$, and we put
\[
  \varepsilon_0=\tfrac1{1000},\qquad K_0=K_{\varepsilon_0},\qquad\kappa(h)=2K_0\,2^{-h/2}.
\]

\begin{proposition}[Cuts of size at most three]\label{prop:threecut}
There is a constant $r_3$ with the following property. Let $G$ be a connected cubic multigraph
of rank $r\ge r_3$ that has a cyclic cut of size at most three. Then
\[
  c(G)\;\le\;2^{r-1}+5K_0^{2}\,2^{r/2}.
\]
\end{proposition}

\begin{proof}
Let $d\le3$ be the least size of a cyclic cut, and let $\delta X$ be a cyclic cut of size $d$. By
the remark at the start of this section, $G$ has no cut of size less than $d$. So $G[X]$ and
$G-X$ are connected, and no vertex is on two edges of $\delta X$ if $d\ge2$: otherwise we could
move that vertex to the other side and get a smaller cut. Let $T_1$ and $T_2$ be the two sides,
with ranks $b_1\le b_2$. Then $b_1\ge1$ and $r=b_1+b_2+d-1$.

If $d=1$, then $c(G)=c(T_1)+c(T_2)\le(2^{b_1}-1)+(2^{b_2}-1)\le2^{r-1}$, because $b_1+b_2=r$.

Let $d\in\{2,3\}$. Then $G$ has no loop, so $T_1$ and $T_2$ are $d$-poles, and the cut edges
match their ports. A cycle
that crosses the cut uses exactly two cut edges, and it consists of these and of one path in each
side. For a pair $j$ of cut edges let $x_j\,2^{b_1}$ and $a_j\,2^{b_2}$ be the numbers of paths
between the corresponding ports of $T_1$ and of $T_2$. Then $0\le x_j,a_j\le1$
by~\eqref{eq:joinbound}, and
\begin{equation}\label{eq:glue23}
  c(G)\;=\;c(T_1)+c(T_2)+2^{b_1+b_2}\sum_jx_ja_j,
  \qquad 2^{r-1}=2^{d-2}\,2^{b_1+b_2}.
\end{equation}
There is one pair for $d=2$, and there are three pairs for $d=3$. Let $r$ be large. Then
$b_2\ge(r-2)/2$ is large, and $\kappa(b_2)\le\frac1{1000}$. We have $c(T_1)<2^{b_1}$.

\step{Case A: $a_j\le0.501$ for all $j$.} By Corollary~\ref{cor:cycles},
$c(T_2)\le\frac9{16}2^{b_2}\le\frac9{32}2^{b_1+b_2}$. So $c(G)\,2^{-b_1-b_2}$ is less than
$2^{-b_2}+\frac9{32}+0.501$ for $d=2$, and less than $2^{-b_2}+\frac9{32}+1.503$ for $d=3$. In
both cases $c(G)<2^{r-1}$.

\step{Case B: $a_j>0.501$ for some $j$.} Then $c(T_2)\le K_0\,2^{b_2/2}$ by
Lemma~\ref{lem:gadgets}(a). So $c(T_1)+c(T_2)<2^{b_1}+K_0\,2^{b_2/2}\le\frac1{1000}2^{b_1+b_2}$.

Let $d=2$. If $x_1<0.501$, then $c(G)<0.502\cdot2^{b_1+b_2}<2^{r-1}$ by~\eqref{eq:glue23}.
Otherwise $c(T_1)\le K_0\,2^{b_1/2}$ by Lemma~\ref{lem:gadgets}(a), and
\[
  c(G)\;\le\;2^{r-1}+K_0\bigl(2^{b_1/2}+2^{b_2/2}\bigr)\;\le\;2^{r-1}+2K_0\,2^{r/2}.
\]

Let $d=3$, and number the pairs so that $a_1\ge a_2\ge a_3$. Any two pairs of ports of a $3$-pole
have one port in common. If $a_1+a_2<1.002$, then $\sum_jx_ja_j<1.503$ and $c(G)<2^{r-1}$.
Otherwise $a_3\le\kappa(b_2)$ by Lemma~\ref{lem:gadgets}(b). If $x_1+x_2<1.002$, then
$\sum_jx_ja_j<1.003$ and $c(G)<2^{r-1}$. Otherwise $x_3\le\kappa(b_1)$ and
$c(T_1)\le K_0\,2^{b_1/2}$ by Lemma~\ref{lem:gadgets}(b) and~(a). Then~\eqref{eq:glue23} gives
\[
  c(G)\;\le\;2K_0\,2^{b_2/2}+2^{b_1+b_2}\bigl(2+\kappa(b_1)\kappa(b_2)\bigr)
  \;=\;2^{r-1}+2K_0\,2^{b_2/2}+4K_0^{2}\,2^{(b_1+b_2)/2}
  \;\le\;2^{r-1}+5K_0^{2}\,2^{r/2}.\qedhere
\]
\end{proof}

\subsection{A pole with four ports}

Let $Y$ be a $4$-pole of rank $h$ with ports $1,2,3,4$. A \emph{pairing} is one of the three
partitions $12|34$, $13|24$, $14|23$ of the ports into two pairs; the two pairs of a pairing are
\emph{opposite}. For a pair $ij$ and a pairing $\alpha=ij|kl$ put
\[
  a_{ij}=2^{-h}P_{ij}(Y),\qquad u_\alpha=2^{-h}f_{ij|kl}(Y),\qquad
  S=\sum_{ij}a_{ij},\qquad F=\sum_\alpha u_\alpha .
\]
A pair of disjoint paths that joins the four ports is a join for the set of all four ports, and
there are $2^{h}$ such joins. So $a_{ij}\le1$ and $F\le1$.

\begin{figure}[ht]
\centering
\begin{tikzpicture}[x=1cm,y=1cm]
  \foreach \ox/\lab in {0/a,6.5/c}{
    \draw[region] (\ox,0) ellipse [x radius=1.0, y radius=1.5];
    \node[text=accent] at (\ox-0.4,0) {$Y$};
    \foreach \i/\y in {1/1.05,2/0.35,3/-0.35,4/-1.05}{
      \coordinate (\lab p\i) at (\ox+0.62,\y);
      \node[note, text=accent] at (\ox+0.3,\y) {\scriptsize\i};}
  }
  % the cycle C
  \foreach \i/\y in {1/1.05,2/0.35,3/-0.35,4/-1.05} \coordinate (c\i) at (2.1,\y);
  \draw[edge] (ap1)--(c1) (ap2)--(c2) (ap3)--(c3) (ap4)--(c4);
  \draw[cool] (c1)--(c2)--(c3)--(c4);
  \draw[cool] (c4) .. controls (3.1,-0.9) and (3.1,0.9) .. (c1);
  \foreach \i in {1,2,3,4} \node[vtx] at (c\i) {};
  \node[note, text=accent] at (2.45,0) {$C$};
  % the cap for 14|23
  \coordinate (x2) at (9.0,0);
  \coordinate (y2) at (8.0,0);
  \draw[hot] (cp1)--(x2)--(cp4) (cp2)--(y2)--(cp3) (x2)--(y2);
  \node[vtx] at (x2) {};
  \node[vtx] at (y2) {};
  \foreach \lab in {a,c} \foreach \i in {1,2,3,4} \node[port] at (\lab p\i) {};
  \node[note] at (1.1,-1.95) {the graph $G$};
  \node[note] at (7.6,-1.95) {the cap for $14|23$};
\end{tikzpicture}
\caption{A $4$-pole $Y$. Left: the graph $G$ of Section~\ref{sec:fourcycle}, in which a cycle
$C$ of length four is joined to the four ports. Right: the cap of Lemma~\ref{lem:cap} for the
pairing $14|23$. The number $a_{ij}$ counts the paths between the ports $i$ and $j$ of $Y$, and
$u_\alpha$ counts the pairs of disjoint paths for a pairing $\alpha$ of the ports, both as a
proportion of the largest possible number.}
\label{fig:glue}
\end{figure}

\begin{lemma}[Cap]\label{lem:cap}
For every $\eta>0$ there is $h_1(\eta)$ with the following property. Let $Y$ be a $4$-pole of
rank $h\ge h_1(\eta)$, and put $z=2^{-h}c(Y)$. Then, for every pairing $\alpha$,
\[
  z+S+F-u_\alpha\;\le\;4+\eta .
\]
\end{lemma}

\begin{proof}
Let $\alpha=ij|kl$. Add two new vertices $x$ and $y$ and an edge between them, join $x$ to the
ports $i$ and $j$, and join $y$ to the ports $k$ and $l$ (Figure~\ref{fig:glue}). The new
multigraph $\hat Y$ has two more vertices and five more edges than $Y$, so its rank is $h+3$. The
new edges form a tree, and this tree joins any two ports by a path. A path of $Y$ between two
ports and the path of the tree between its ends form a cycle. A pair of disjoint paths of $Y$
whose pairing is not $\alpha$ and the four edges of the tree at the ports form a cycle. With the
cycles of $Y$, these are $c(Y)+(S+F-u_\alpha)2^{h}$ distinct cycles of $\hat Y$. By
Corollary~\ref{cor:cycles}, $\hat Y$ has at most $(\half+\frac18\eta)\,2^{h+3}=(4+\eta)\,2^{h}$
cycles, when $h$ is large.
\end{proof}

By Lemma~\ref{lem:cap}, the number $S+F-u_\alpha$ is at most $4+o(1)$ for every pairing
$\alpha$. The next lemma describes the poles in which $S+F$ is at least $4-o(1)$. They have one
of two shapes (Figure~\ref{fig:shapes}).

\begin{figure}[ht]
\centering
\begin{tikzpicture}[x=1cm,y=1cm]
  \begin{scope}
    \draw[cool] (0,1)--(0,0) (0,1)--(1.7,1) (0,1)--(1.7,0);
    \foreach \p in {(0,1),(0,0),(1.7,1),(1.7,0)} \node[port] at \p {};
    \node[note, text=accent] at (-0.3,1) {\scriptsize 1};
    \node[note, text=accent] at (-0.3,0) {\scriptsize 2};
    \node[note, text=accent] at (2.0,1) {\scriptsize 3};
    \node[note, text=accent] at (2.0,0) {\scriptsize 4};
    \node[note] at (0.85,-0.6) {three pairs with a common port};
  \end{scope}
  \begin{scope}[xshift=7.0cm]
    \draw[cool] (0,1)--(1.7,1) (0,1)--(1.7,0) (0,0)--(1.7,1) (0,0)--(1.7,0);
    \foreach \p in {(0,1),(0,0),(1.7,1),(1.7,0)} \node[port] at \p {};
    \node[note, text=accent] at (-0.3,1) {\scriptsize 1};
    \node[note, text=accent] at (-0.3,0) {\scriptsize 2};
    \node[note, text=accent] at (2.0,1) {\scriptsize 3};
    \node[note, text=accent] at (2.0,0) {\scriptsize 4};
    \node[note] at (0.85,-0.6) {the four pairs between two pairs of ports};
  \end{scope}
\end{tikzpicture}
\caption{The two shapes of Lemma~\ref{lem:rigid}. A line joins two ports whose pair is full.}
\label{fig:shapes}
\end{figure}

\begin{lemma}[Rigidity]\label{lem:rigid}
There are numbers $\rho_0\in(0,1)$ and $h_0$ with the following property. Let $Y$ be a $4$-pole
of rank $h\ge h_0$ with $S+F\ge4-\rho_0$. Then three pairs with a common port are full, or the
four pairs between two disjoint pairs of ports are full.
\end{lemma}

\begin{proof}
Suppose that this is false. Then there are $4$-poles $Y_1,Y_2,\dots$ whose ranks $h_n$ tend to
infinity, with $S(Y_n)+F(Y_n)\ge4-\frac1n$, such that no $Y_n$ has three full pairs with a common
port or four full pairs between two disjoint pairs of ports. The numbers $a_{ij}(Y_n)$ and
$u_\alpha(Y_n)$ lie in $[0,1]$. We take a subsequence in which each of these nine sequences
converges, and we write $a_{ij}$, $u_\alpha$, $S$ and $F$ for the limits. Then $F\le1$ and
$S+F\ge4$.

\step{Three implications.} For distinct ports $i,j,k,l$,
\begin{align}
  &a_{ik}+a_{jk}>1\ \Longrightarrow\ a_{ij}=0;\label{eq:T1}\\
  &a_{ik}+a_{il}+a_{jk}+a_{jl}>2\ \Longrightarrow\ u_{ik|jl}+u_{il|jk}=0;\label{eq:T2}\\
  &a_{kl}>0\ \Longrightarrow\ a_{ij}+F-u_{ij|kl}\le1.\label{eq:T3}
\end{align}
Indeed, let $a_{ik}+a_{jk}>1+2\varepsilon$ with $\varepsilon>0$. For large $n$,
Lemma~\ref{lem:gadgets}(b) gives $a_{ij}(Y_n)\le2K_\varepsilon\,2^{-h_n/2}$, and this tends to
$0$. In the same way~\eqref{eq:T2} follows from Lemma~\ref{lem:gadgets}(c). If
$a_{ij}+F-u_{ij|kl}>1$, then Lemma~\ref{lem:gadgets}(d) gives $a_{kl}=0$; this
is~\eqref{eq:T3}.

\step{The limit.} For a partition of the ports into two pairs, the \emph{cross sum} is the sum of
$a_{ij}$ over the four pairs $ij$ that are not pairs of the partition. Each pair occurs in two of
the three cross sums.

Suppose that some cross sum is larger than $2$, say the cross sum $X$ of $12|34$. Then
$a_{13}+a_{23}>1$ or $a_{14}+a_{24}>1$, so $a_{12}=0$ by~\eqref{eq:T1}. In the same way
$a_{34}=0$. By~\eqref{eq:T2}, $u_{13|24}=u_{14|23}=0$. Put $u=u_{12|34}$. Then $X+u=S+F\ge4$.
If one of the four numbers in $X$ is $0$, say $a_{24}$, then the other three and $u$ are equal
to $1$. This contradicts~\eqref{eq:T3} for $ij=14$, because $a_{23}>0$. So the four numbers are
positive, and~\eqref{eq:T3} gives $a_{ij}+u\le1$ for each of them. Hence $X+4u\le4$. With
$X+u\ge4$ this gives $u=0$ and $X=4$, so the four numbers are equal to $1$.

Suppose that no cross sum is larger than $2$. The three cross sums add up to $2S$, so $S\le3$.
Hence $S=3$ and $F=1$, every cross sum is $2$, and two opposite numbers add up to $1$. If all six
numbers are positive, then~\eqref{eq:T3} gives $u_\alpha\ge\max\{a_{ij},a_{kl}\}\ge\half$ for
every pairing $\alpha=ij|kl$, and this contradicts $F=1$. So one number is $0$, say $a_{34}$, and
$a_{12}=1$. If $a_{13}>0$, then $a_{23}=0$ by~\eqref{eq:T1}, so $a_{14}=1$; then $a_{24}=0$
by~\eqref{eq:T1}, so $a_{13}=1$. If $a_{13}=0$, then $a_{24}=1$; then $a_{14}=0$
by~\eqref{eq:T1}, so $a_{23}=1$. In both cases three numbers with a common port are equal
to $1$.

\step{The end.} So three numbers $a_{ij}$ with a common port, or the four numbers between two
disjoint pairs of ports, are equal to $1$. For large $n$ the corresponding numbers
$a_{ij}(Y_n)$ are larger than $1-\gamma_0$, and these pairs are full in $Y_n$ by
Theorem~\ref{thm:gap}. This is a contradiction.
\end{proof}

\needroom{14\baselineskip}
\subsection{Cycles of length four}\label{sec:fourcycle}

\begin{idea}{outline of the proof of Theorem~\ref{thm:fourcut}}
Let $C$ be a cycle of length four, and let $Y$ be the rest of the graph, of rank $k$. A cycle
that meets $C$ and $Y$ consists of paths of $Y$ and of paths of $C$. In units of $2^{k}$, the
number of these cycles is the sum of the numbers $S+F-u_\alpha$ for the two pairings $\alpha$
that pair the ports at neighbours on $C$. By the cap of Lemma~\ref{lem:cap}, one of these
numbers and the proportion of cycles of $Y$ are together at most $4+o(1)$, and $c(G)\ge2^{r-1}$
requires the sum $8$. So $S+F\ge4-o(1)$. By Lemma~\ref{lem:rigid} there are two possible shapes.
Three full pairs at one port give a cyclic cut of size three. If the four pairs between two
pairs of ports are full, then $Y$ is a chain of tiles, and $G$ is a ring.
\end{idea}

\begin{figure}[ht]
\centering
\begin{tikzpicture}[x=1cm,y=1cm]
  \coordinate (s) at (0,0);
  \coordinate (t) at (9.6,0);
  \foreach \n/\x in {a/1.2,b/2.4,c/3.8,d/5.2,e/6.4,f/7.8}{
    \coordinate (\n u) at (\x,0.65);
    \coordinate (\n l) at (\x,-0.65);}
  \draw[region] (-0.4,-1.0) rectangle (2.85,1.0);
  \draw[edge] (s)--(au) (s)--(al);
  \draw[edge] (au)--(bu) (al)--(bl) (au)--(bl) (al)--(bu);
  \draw[hot] (bu)--(cu) (bl)--(cl);
  \draw[edge] (cu)--(cl);
  \draw[edge] (cu)--(du) (cl)--(dl);
  \draw[edge] (du)--(eu) (dl)--(el) (du)--(el) (dl)--(eu);
  \draw[edge] (eu)--(fu) (el)--(fl);
  \draw[cool] (fu)--(fl) (fu)--(t) (fl)--(t);
  % the edge from p to C
  \draw[hot] (s)--(-1.3,0);
  \draw[accent, line width=0.6pt] (-1.65,0) circle [radius=0.35];
  \node[note, text=accent] at (-1.65,0) {$C$};
  % the ports v and v'
  \coordinate (v) at (4.5,0.65);
  \coordinate (vp) at (7.1,-0.65);
  \draw[edge, dash pattern=on 2pt off 1.6pt] (v)--(4.5,1.2) (vp)--(7.1,-1.2) (t)--(10.4,0);
  \foreach \q in {au,al,bu,bl,cu,cl,du,dl,eu,el,fu,fl} \node[vtx] at (\q) {};
  \foreach \q in {s,t,v,vp} \node[port] at (\q) {};
  \node[note, text=accent] at (0,0.35) {$p$};
  \node[note, text=accent] at (9.75,0.35) {$q$};
  \node[note, text=accent] at (4.5,0.33) {$v$};
  \node[note, text=accent] at (7.1,-0.33) {$v'$};
  \node[note, text=accent] at (1.2,1.25) {the set $X$};
  \node[note, text=accent] at (8.45,0) {$Z$};
  \node[note, anchor=west] at (4.6,1.2) {to $C$};
  \node[note, anchor=west] at (7.2,-1.2) {to $C$};
  \node[note, anchor=west] at (10.45,0) {to $C$};
  \node[note, text=warm] at (3.1,-1.3) {three edges leave $X$};
\end{tikzpicture}
\caption{Case~1 in the proof of Theorem~\ref{thm:fourcut}. Without the ports $v$ and $v'$, the
pole $Y$ is a capped ladder with the caps $p$ and $q$. The ports $v$ and $v'$ are not on an
edge inside the set $X$, because a path from $p$ to such a port avoids the cycle $Z$.}
\label{fig:endpiece}
\end{figure}

\begin{theorem}[A cycle of length four]\label{thm:fourcut}
There is a constant $r_4$ with the following property. Let $G$ be a connected cubic multigraph
of rank $r\ge r_4$ that has no cyclic cut of size at most three and has a cycle of length four.
Then $G$ is a ring, or $c(G)<2^{r-1}$.
\end{theorem}

\begin{proof}
Let $C=c_1c_2c_3c_4$ be a cycle of length four, and put $Y=G-V(C)$ (Figure~\ref{fig:glue}). By
the remark at the start of the section, $G$ has no loop and no cut of size at most two, and
every cut of size three has a single vertex on one side. The graph has $2r-2$ vertices, so a set
of at most five vertices and its complement both have more than one vertex.

\step{The rest is a $4$-pole without bridge.} No edge outside $C$ joins two vertices of $C$:
such an edge and a path of $C$ form a cycle with at most three vertices, and the cut around
these vertices has size at most three. So exactly four edges leave $V(C)$, one at each vertex.
Let $c_i$ be joined to the vertex $i$ of $Y$. The vertices $1,2,3,4$ are distinct: if two of
them are one vertex $v$, the cut around $V(C)\cup\{v\}$ has size at most three. The multigraph
$Y$ is connected: otherwise some component of $Y$ is joined to $C$ by at most two edges, and
these form a cut. So $Y$ is a $4$-pole with the ports $1,2,3,4$. Let $k$ be its rank. Then
$r=k+4$.

Let $e$ be a bridge of $Y$, and let $A$ and $B$ be the vertex sets of the two components of
$Y-e$. In $G[A]$ every vertex has degree at least two, except the end of $e$, which has degree at
least one. A forest with an edge has two vertices of degree one, so $G[A]$ contains a cycle. The
same holds for $G[B]$. Each of the four edges between $C$ and $Y$ is in exactly one of the cuts
$\delta A$ and $\delta B$, and $e$ is in both. So $|\delta A|+|\delta B|=6$, and one of them is a
cyclic cut of size at most three. Hence $Y$ has no bridge.

\step{The count.} We use the notation of this section for $Y$, with $h=k$, and we put
$z=2^{-k}c(Y)$. A cycle of $G$ that is neither $C$ nor a cycle of $Y$ uses two or four of the
edges between $C$ and $Y$. With two, at $c_i$ and $c_j$, it consists of a path of $Y$ between the
ports $i$ and $j$ and of one of the two paths of $C$ between $c_i$ and $c_j$. With four, it
consists of a pair of disjoint paths of $Y$ and of two disjoint edges of $C$. The edges $c_1c_2$
and $c_3c_4$ close the pair to one cycle if and only if its pairing is not $12|34$. The edges
$c_1c_4$ and $c_2c_3$ do this if and only if its pairing is not $14|23$. Hence
\begin{equation}\label{eq:four}
  \frac{c(G)}{2^{k}}\;=\;2^{-k}+z+\bigl(S+F-u_{12|34}\bigr)+\bigl(S+F-u_{14|23}\bigr).
\end{equation}
Let $\rho_0$ and $h_0$ be as in Lemma~\ref{lem:rigid}, and let $r$ be so large that
$k\ge\max\{h_0,\,h_1(\rho_0/2),\,4\}$ and $2^{-k}\le\rho_0/2$. Suppose that
$c(G)\ge2^{r-1}=8\cdot2^{k}$. We show that $G$ is a ring. By Lemma~\ref{lem:cap} for the pairing
$14|23$, the term $z$ and the second bracket of~\eqref{eq:four} are together at most
$4+\rho_0/2$. So $S+F\ge S+F-u_{12|34}\ge4-\rho_0$, and Lemma~\ref{lem:rigid} applies to $Y$. It
gives two cases.

\step{Case 1: three pairs with a common port are full.} Let $p$ be the common port, and let $q$,
$v$ and $v'$ be the other ports. Replace $v$ and $v'$, each with its two edges, by one edge. By
Lemma~\ref{lem:reduction}(a) the result $Y^{\circ}$ is a $2$-pole of rank $k$ with the ports
$p,q$, and this pair is full. It has no bridge, because $Y$ has none. By
Lemma~\ref{lem:ladders}, $Y^{\circ}$ is a capped ladder with the caps $p$ and $q$. Its tiles
have total weight $k-1\ge3$, so there are at least two tiles. Let $X$ consist of $p$ and of the
vertices of the first tile (Figure~\ref{fig:endpiece}), and let $Z$ be a cycle of $Y^{\circ}$ in the last tile and the cap
$q$. Then $Z$ avoids $X$, and $Y^{\circ}[X]$ is connected and contains a cycle.

Suppose that the edge of $Y^{\circ}$ that replaces $v$ has both ends in $X$. Then $Y$ has a path
from $p$ to $v$ whose other vertices are in $X\cup\{v'\}$. It is disjoint from the cycle of $Y$
that corresponds to $Z$. This contradicts Lemma~\ref{lem:joins} for the full pair $p,v$. The
same holds for $v'$. So the edges of $Y^{\circ}$ with both ends in $X$ are edges of $Y$, and
$G[X]$ contains a cycle. Three edges of $G$ leave $X$: the edge from $p$ to $C$, and one edge of
$Y$ for each of the two edges of $Y^{\circ}$ that leave $X$. The cycle $C$ is outside $X$. So
$\delta X$ is a cyclic cut of size three, and this case does not occur.

\step{Case 2: the four pairs between two disjoint pairs of ports are full.} The pole $Y$ has no
bridge. By Corollary~\ref{cor:fourports} it is a chain of tiles; one of the two pairs of ports
is the left pair of its first tile, and the other is the right pair of its last tile. If the
pairs are $1,3$ and $2,4$, then $C$ is a square with the pairs $c_1,c_3$ and $c_2,c_4$. In the
other case let the pairs be $1,2$ and $3,4$, by symmetry. Then the edges $c_1c_2$ and $c_3c_4$
are single rungs, and the edges $c_2c_3$ and $c_1c_4$ are the interface between them. In both
cases the edges between $C$ and $Y$ are two interfaces. So $G$ is a ring.
\end{proof}

\section{Large girth outside a small set}\label{sec:half}

Recall that $q(G,P)$ is the probability that $\xi\cap E(G-P)$ is a forest.

\begin{theorem}\label{thm:half}
There are integers $D_*\ge7$ and $r_5$ with the following property. Let $G$ be a connected cubic
multigraph of rank $r\ge r_5$. Let $P$ be a nonempty set of at most $3\log r$ vertices such that
$G[P]$ is connected. If $G-P$ has girth at least $D_*$, then $q(G,P)<\frac25$.
\end{theorem}

\begin{idea}{outline of the proof}
Suppose that $q(G,P)\ge\frac25$. We find two disjoint parts of $G-P$ that are large, have small
boundaries and have large girth. Each of them has a forest proportion of at most $\frac35$, for
every choice of the edges outside it. This is a contradiction, because $(\frac35)^{2}<\frac25$. To
find the parts we use a packing bound for short cycles (Section~\ref{sec:packing}), a second
moment and a chain of nested separators (Section~\ref{sec:chain}).
\end{idea}

\subsection{Pieces with a small boundary}\label{sec:affine}

Let $N$ be a pole, and let $T$ be a set of ports of even size. A \emph{$T$-join} of $N$ is a set
of edges in which the vertices of $T$ have odd degree and all other vertices have even degree.
There are $2^{\beta(N)}$ of them. The \emph{forest proportion} of $T$ is the proportion of the
$T$-joins that are forests. For $T=\emptyset$ the $T$-joins are the even subgraphs, only the empty
one is a forest, and the forest proportion is $2^{-\beta(N)}$. Put
\[
  R_{*}=R_{1/20},\qquad B=2R_{*}+2^{100}+1,
\]
with $R_{1/20}$ as in Proposition~\ref{prop:input}. Then $\Theta_{1/20}(k)+k\le B(k+1)^{2}$ for all
$k\ge1$.

\begin{lemma}[A marker in a pole]\label{lem:marker}
Let $N$ be a pole with $k\ge2$ ports, and let $o$ be a port. Call the set of the vertices at
distance $j$ from $o$ the \emph{layer} $j$. Suppose that some set $T$ of ports has forest
proportion at least $\frac35$. Then each layer has at most $B(k+1)^{2}$ vertices.
\end{lemma}

\begin{proof}
Take a breadth-first tree of $N$ from $o$, and let $P$ be the union of its paths from $o$ to the
other ports. Then $N[P]$ is connected, $P$ contains all vertices of degree two, and $P$ has at most
$k$ vertices in each layer, because a path of the tree from $o$ has one vertex in each layer that
it meets. The tree paths inside $P$ contain a $T$-join $F_0$. The map $J\mapsto J+F_0$ is a
bijection from the $T$-joins of $N$ to the even subgraphs of $N$, and $J+F_0$ has the same edges in
$N-P$ as $J$. So if $J$ is a forest, then $J+F_0$ is a forest in $N-P$. Hence
$q(N,P)\ge\frac35>\half+\frac1{20}$. The distance from $o$ is a level function. By
Lemma~\ref{lem:width} with $\delta=\frac1{20}$, each layer has fewer than $\Theta_{1/20}(k)$
vertices outside $P$, and it has at most $k$ vertices in $P$.
\end{proof}

\begin{lemma}[Girth and rank of a pole]\label{lem:moore}
Let $N$ be a pole of rank $h\ge2$ with $k$ ports and girth at least $D$. Put
$t=\lceil\log(8h)\rceil$. Then
\[
  D\;\le\;4t\Bigl(1+\frac{k}{3(h-1)}\Bigr).
\]
\end{lemma}

\begin{proof}
Replace every vertex of degree two and its two edges by one edge, and give each edge of the result
the length of the path of $N$ that it replaces. The result $N_0$ is a connected cubic multigraph
with $n_0=2h-2$ vertices and $m_0=3h-3$ edges, and the lengths of its edges add up to $m_0+k$. A
cycle of $N_0$ comes from a cycle of $N$ of the same length.

There are $2m_0\,2^{t-1}$ non-backtracking walks with $t$ edges in $N_0$. Choose one of them
uniformly. Then each of its steps is uniform over the $2m_0$ directed edges. So the expected length
of the walk is $t(1+k/m_0)$. At least half of the walks have length at most twice this value. Their
number is at least
\[
  m_0\,2^{t-1}\;\ge\;12h(h-1)\;>\;n_0^{2}.
\]
If one of these short walks visits a vertex twice, it contains a cycle that is not longer than the
walk. Otherwise all of them are paths, and two of them have the same first vertex and the same last
vertex. Their union contains a cycle of length at most the sum of their lengths. In both cases $N$
has a cycle of length at most $4t(1+k/m_0)$.
\end{proof}

\begin{lemma}[Forest proportions at large girth]\label{lem:affine}
There are absolute constants $D_{*}\ge7$ and $C_0$ with the following property. Let $N$ be a pole
of rank $h$ with $k$ ports and girth at least $D$. Suppose that
\begin{enumerate}
\item[\textup{(a)}] $h\ge256$, $k\le4D$ and $D\ge D_{*}$, or
\item[\textup{(b)}] $h\ge\Lambda\ge4$, $k\le2\Lambda$ and $D\ge C_0\log(\Lambda+2)$ for some
number $\Lambda$.
\end{enumerate}
Then every set of ports of even size has forest proportion at most $\frac35$.
\end{lemma}

\begin{proof}
If $k\le1$, the only set is the empty one, with forest proportion $2^{-h}$. Let $k\ge2$, and put
$t=\lceil\log(8h)\rceil\le\log h+4$.

\step{The rank is exponential in the girth.} In case~(a), $h-1>\frac{32}3t$: this holds at
$h=256$, where $t=11$, and the ratio $(h-1)/(\log h+4)$ increases with $h$. So
Lemma~\ref{lem:moore} gives $D\le4t+16tD/(3(h-1))<4t+D/2$. In case~(b),
$k/(3(h-1))\le2h/(3(h-1))<1$, and Lemma~\ref{lem:moore} gives $D<8t$. In both cases
$D<8t\le8\log h+32$, so
\begin{equation}\label{eq:rankgirth}
  h\;\ge\;2^{D/8-4},\qquad\text{and }N\text{ has }n\ge h\text{ vertices}.
\end{equation}

\step{A ball without ports.} Put $s=\lfloor D/16\rfloor$. A ball of radius $s$ has at most
$3\cdot2^{s}$ vertices. So at most $3k\,2^{s}$ vertices are within distance $s$ of a port. In
case~(a) this is at most $12D\,2^{D/16}$, and in case~(b) at most $6\Lambda\,2^{D/16}$. Choose
$D_{*}$ and $C_0$ so large that this number is less than $2^{D/8-4}\le n$ in each case. Then some
vertex $w$ has distance more than $s$ from every port. All vertices within distance $s$ of $w$
have degree three, and $2s<D$. So the ball of radius $s$ around $w$ is a tree in which every
vertex at distance less than $s$ from $w$ has degree three, and it has at least
$2^{s}\ge2^{D/16-1}$ vertices. Fix a port $o$. The distances from $o$ to two vertices of this ball
differ by at most $2s$. So the ball meets at most $2s+1\le D$ layers with respect to $o$, and some
layer has at least $2^{D/16-1}/D$ vertices.

\step{Conclusion.} Suppose that some set of ports has forest proportion more than $\frac35$. By
Lemma~\ref{lem:marker} every layer has at most $B(k+1)^{2}$ vertices. This is at most
$B(4D+1)^{2}$ in case~(a) and at most $B(2\Lambda+1)^{2}$ in case~(b). In case~(a) choose $D_{*}$ so
large that $2^{D/16-1}/D>B(4D+1)^{2}$ for $D\ge D_{*}$. In case~(b) the function $2^{D/16}/D$
increases for $D\ge32$, and at $D=C_0\log(\Lambda+2)$ its ratio to $(\Lambda+2)^{2}$ is
$(\Lambda+2)^{C_0/16-2}/(C_0\log(\Lambda+2))$. This tends to infinity with $C_0$, uniformly in
$\Lambda\ge4$. So a large $C_0$ gives $2^{D/16-1}/D>B(2\Lambda+1)^{2}$. In both cases we have a
contradiction.
\end{proof}

\paragraph{Products of pieces.}

\begin{lemma}[Product]\label{lem:product}
Let $G$ be a cubic multigraph, let $P$ be a set of vertices, and let
$U_1,\ldots,U_m\subseteq V(G)\sm P$ be pairwise disjoint sets such that each $G[U_i]$ is a pole.
Suppose every even set of ports in each pole has forest proportion at most $3/5$.
Then $q(G,P)\le(3/5)^m$.
\end{lemma}

\begin{proof}
Let $E'$ consist of the edges that lie in none of the induced poles.
Fix a restriction $x\subseteq E'$ that extends to an even subgraph of $G$.
For each $i$, let $T_i$ be the vertices of $U_i$ incident with an odd number of edges of $x$.
The parity equations say that $\xi\cap E(G[U_i])$ is a $T_i$-join.
These equations involve disjoint internal edge sets, so the joins are independent and uniform,
conditional on $\xi\cap E'=x$.
A vertex of degree three in a pole has no incident edge in $E'$; thus each $T_i$ is an even set of ports.
If $\xi\cap E(G-P)$ is a forest, its restriction to every pole is a forest.
The conditional probability is therefore at most $(3/5)^m$. Average over $x$.
Edges between the poles belong to $E'$, so adjacent pieces cause no exception.
\end{proof}

\begin{lemma}[The core of a piece]\label{lem:core}
Let $G$ be a cubic multigraph, and let $X$ be a set of vertices such that $G[X]$ is connected, has
no loop, and has rank $h\ge2$. Then there is a set $U\subseteq X$ such that $G[U]$ is a pole of rank $h$ with at
most $|\delta X|$ ports, and the girth of $G[U]$ is at least the girth of $G[X]$.
\end{lemma}

\begin{proof}
Delete vertices of degree one from $G[X]$ as long as there are any. The result $G[U]$ is connected,
has the same cycles as $G[X]$ and all its degrees are two or three. When a vertex of degree one is
deleted, two edges that left the set no longer do, and one edge inside the set now leaves it. So
$|\delta U|\le|\delta X|$. A vertex of degree two in $G[U]$ has exactly one edge end outside
$G[U]$, and a vertex of degree three has none. So $G[U]$ has $|\delta U|$ ports.
\end{proof}

\begin{proposition}[Pieces with small boundaries]\label{prop:blocks}
Let $G$ be a connected cubic multigraph and let $P$ be a set of vertices. Let
$X_1,\ldots,X_m\subseteq V(G)\sm P$ be pairwise disjoint, nonempty and proper sets of vertices. Suppose that one of
the following holds for every $i$.
\begin{enumerate}
\item[\textup{(a)}] $G[X_i]$ is connected of rank at least $256$ and girth at least $D\ge D_{*}$,
and $|\delta X_i|\le4D$.
\item[\textup{(b)}] $|X_i|\ge8\Lambda^{2}$, $|\delta X_i|\le2\Lambda$, and $G[X_i]$ has girth at
least $D$, where $\Lambda\ge4$ and $D\ge C_0\log(\Lambda+2)$.
\end{enumerate}
Then $q(G,P)\le(3/5)^m$.
\end{proposition}

\begin{proof}
It suffices to find sets $U_i\subseteq X_i$ as in Lemma~\ref{lem:product}. In case~(a),
Lemma~\ref{lem:core} gives a pole $G[U_i]$ of rank at least $256$ with at most $4D$ ports and girth
at least $D$, and Lemma~\ref{lem:affine}(a) applies.

In case~(b) write $n=|X_i|$, $d=|\delta X_i|$ and $c=\comp(G[X_i])$. Every component of $G[X_i]$
has an edge that leaves $X_i$, so $1\le c\le d\le2\Lambda$. Count the edge ends:
$3n=2|E(G[X_i])|+d$. So $\beta(G[X_i])=(n-d)/2+c$, and some component $Y$ of $G[X_i]$ has rank at
least
\[
  \frac{n-d}{2c}+1\;\ge\;\frac{8\Lambda^{2}-2\Lambda}{4\Lambda}+1\;>\;\Lambda .
\]
Its boundary is a part of $\delta X_i$. Lemma~\ref{lem:core} for $Y$ gives a pole $G[U_i]$ of rank
at least $\Lambda$ with at most $2\Lambda$ ports and girth at least $D$, and
Lemma~\ref{lem:affine}(b) applies.
\end{proof}

\subsection{Packing of short cycles}\label{sec:packing}

\begin{idea}{the purpose of the packing bound}
The second-moment argument requires a graph of girth about $\log r$. So we delete a maximal
family of disjoint short cycles, and we need a bound for the size of this family.
Lemma~\ref{lem:smallcuts} gives such a bound only when the proportion of forests exceeds one half
by a fixed margin. Here we prove the bound with Proposition~\ref{prop:blocks}. Many disjoint
cyclic sets with small boundaries have independent local events, or many of them lie in one
string (Lemma~\ref{lem:strings}), and a long string contains two pieces to which
Proposition~\ref{prop:blocks} applies.
\end{idea}

Throughout, $h_*=256$.

\paragraph{Vertices that share a parity.}

Let $Q$ be a connected multigraph without loops. For an edge $e$, the map
$\xi\mapsto[e\in\xi]$ is a linear functional on $\Z(Q)$. For a vertex $v$ let $L_v$ be the space
of functionals that is spanned by the edges at $v$. For a set $A$ of vertices put
$U_A=\sum_{a\in A}L_a$. We use one standard fact: if a set $F$ of edges meets every even subgraph of
$Q$ in an even number of edges, then $F=\delta Y$ for some set $Y$ of vertices. (Fix a spanning
tree and a root. Let $Y$ be the set of the vertices whose tree path from the root has an odd number
of edges in $F$. Then $F$ and $\delta Y$ contain the same tree edges, and the fundamental cycles
show that they contain the same edges outside the tree.)

\begin{lemma}[Strings]\label{lem:strings}
Let $Q$ be a connected multigraph without loops, and let $A$ be a nonempty proper set of its
vertices. Put $b=|\delta A|$. Call a vertex $v\notin A$ \emph{relevant} if $L_v\cap U_A\ne0$. There
are a set of at most $2b$ relevant vertices, the \emph{exceptions}, and at most $2b$ sequences of
relevant vertices, the \emph{strings}, with these properties. Every relevant vertex that is not an
exception is in a string. If $v_1,\dots,v_m$ is a string and $1\le i<j\le m$ with $j\ge i+2$, there
is a set $M(i,j)$ of vertices such that
\begin{enumerate}
\item[\textup{(i)}] $M(i,j)$ is disjoint from $A\cup\{v_i,v_j\}$, the graph $Q[M(i,j)]$ is
connected, and $v_l\in M(i,j)$ for $i<l<j$;
\item[\textup{(ii)}] every edge that leaves $M(i,j)$ ends in $v_i$ or in $v_j$;
\item[\textup{(iii)}] $M(i,j)\cap M(i',j')=\emptyset$ if $j\le i'$.
\end{enumerate}
\end{lemma}

\begin{proof}
Call the edges of $\delta A$ \emph{terminals}. The end of a terminal outside $A$ is its
\emph{foot}. Every component of $Q-A$ contains a foot, because $Q$ is connected.

\step{Step 1.} Let $v\notin A$ be relevant and not a foot, and let $K$ be the component of $Q-A$
that contains $v$. We show that at least two components of $K-v$ contain feet. Let
$0\ne a\in L_v\cap U_A$. Then $a$ is the sum of the functionals of a set $F$ of edges at $v$, and
also of a set $F'$ of edges with an end in $A$. No edge is in both sets, because $v$ is not a foot.
The set $F\cup F'$ meets every even subgraph in an even number of edges, so $F\cup F'=\delta Y$ for
some $Y$. We may assume that $v\notin Y$. No edge of $K-v$ is in $F\cup F'$. So each component of
$K-v$ lies inside $Y$ or outside $Y$, and $F$ is the set of all edges from $v$ to the components
inside $Y$. For a component $J$ of $K-v$, the cut $\delta V(J)$ consists of the edges from $J$ to
$v$ and the terminals with a foot in $J$. So on $\Z(Q)$ the sum of the edges from $J$ to $v$ equals
the sum of the terminals with a foot in $J$. Hence $a$ is the sum of the terminals with a foot in
$Y$. The sum of all terminals with a foot in $K$ is zero on $\Z(Q)$, because these terminals form
the cut $\delta V(K)$. As $a\ne0$, some component of $K-v$ inside $Y$ contains a foot, and some
component outside $Y$ contains a foot.

\step{Step 2.} In each component $K$ of $Q-A$ fix a spanning tree. For each terminal add a new leaf
at its foot. Let $T_K$ be the union of the paths of this tree between the new leaves, and let $T$
be the union of the trees $T_K$. Every leaf of $T$ is a new leaf, so $T$ has at most $b$ leaves and
at most $b$ vertices of degree at least three. The exceptions are the relevant vertices that are
feet or have degree at least three in $T$. There are at most $2b$ of them.

Let $v$ be a relevant vertex that is not an exception. By Step~1 two components of $K-v$ contain
feet. The path of the tree between the two new leaves at these feet passes through $v$. So $v$ is
a vertex of $T$, and its degree in $T$ is two. A \emph{thread} is a maximal path of $T$ whose inner
vertices have degree two in $T$. There are at most $2b$ threads, and every vertex of degree two of
$T$ is an inner vertex of exactly one thread. The strings are the sequences of the relevant
vertices that are not exceptions on one thread, in the order of the thread.

\step{Step 3.} Let $v$ be a vertex of a string, in the component $K$. The forest $T-v$ has exactly
two components that meet $K$, one on each side of $v$ along its thread. Their vertex sets in $K$ are
the \emph{near side} of $v$, towards the start of the thread, and the \emph{far side} of $v$. Each
side is connected in $K-v$, and all feet in $K$ lie in the two sides. By Step~1 the two sides lie in
different components of $K-v$.

Now let $v_i$ and $v_j$ be as in the lemma. Let $\Pi$ be the part of their thread strictly between
them. It contains $v_{i+1}$, it is connected, and it contains no foot, because a new leaf at a
vertex of $\Pi$ would give that vertex degree three in $T$. Let $M(i,j)$ be the vertex set of the
component of $K-v_i-v_j$ that contains $\Pi$. This gives~(i).

For~(ii), suppose that $M(i,j)$ contains a foot $w$. Then $w$ is a vertex of $T$ outside $\Pi$. So
$w$ lies in the near side of $v_i$ or in the far side of $v_j$; say in the near side of $v_i$. The
path $\Pi$ lies in the far side of $v_i$. But $w$ and $\Pi$ are joined in $K-v_i$ through
$M(i,j)$, and the two sides of $v_i$ lie in different components of $K-v_i$. This is a
contradiction. So $M(i,j)$ contains no foot. An edge that leaves $M(i,j)$ ends in $A$, in $v_i$ or
in $v_j$, and it does not end in $A$.

For~(iii), let $j\le i'$. The set $M(i,j)$ is connected in $K-v_j$ and contains $\Pi$, which lies
in the near side of $v_j$. So $M(i,j)$ lies in the component of $K-v_j$ that contains the near side
of $v_j$. We show that $M(i',j')$ lies in the component of $K-v_j$ that contains the far side of
$v_j$; then the two sets are disjoint. The set $M(i',j')$ is connected in $K-v_{i'}$ and contains
the part $\Pi'$ of the thread between $v_{i'}$ and $v_{j'}$, which lies in the far side of
$v_{i'}$. If $j=i'$, this is the claim. If $j<i'$, then $v_j$ lies in the near side of $v_{i'}$, so
$v_j\notin M(i',j')$. Hence $M(i',j')$ is connected in $K-v_j$, and it contains $\Pi'$, which lies
in the far side of $v_j$.
\end{proof}

\paragraph{A packing bound.}

\begin{proposition}[Packing]\label{prop:packing}
Let $G$ be a connected cubic multigraph. Let $P$ be empty or induce a connected subgraph,
and let $S$ consist of $s$ vertices outside $P$.
Suppose $G-(P\cup S)$ has girth at least $D\ge D_*$ and $q(G,P\cup S)\ge\frac25$.
Let $2\le d\le2D$, and let $X_1,\ldots,X_n$ be pairwise disjoint sets outside $P\cup S$,
each inducing a connected graph with a cycle and satisfying $|\delta X_i|\le d$.
Then $n\le1551\,d(s+1)2^{d}$.
\end{proposition}

\begin{proof}
Contract each $X_i$ to a marked vertex $x_i$, contract each component of
$G-\bigcup_iX_i$ to a vertex, and delete the loops. The resulting graph $Q$ is connected.
Projection maps the uniform distribution on $\Z(G)$ to the uniform distribution on $\Z(Q)$,
as in Lemma~\ref{lem:loc}(b). Conditional on the projected edges, the restrictions to the
induced graphs $G[X_i]$ are independent affine joins.
Let $I_i$ be the event that the projection uses no edge at $x_i$.
On $I_i$ the restriction to $G[X_i]$ is an even subgraph, and its forest probability is at most
$1/2$, since $G[X_i]$ contains a cycle. Hence
\begin{equation}\label{eq:packprod}
 \frac25\le q(G,P\cup S)\le
 \EE\prod_{i=1}^n\left(1-\frac{\mathbf1_{I_i}}2\right).
\end{equation}

Write $L_i=L_{x_i}$ for the local parity space in $Q$.
The sum of its incident edge functionals is zero, so $\dim L_i\le d-1$ and
$\PP(I_i)=2^{-\dim L_i}\ge2^{1-d}$.
If $z_0$ local spaces are zero, then $\frac25\le2^{-z_0}$.
Thus at most one local space is zero.
Keep its vertex in $Q$, but omit this mark from the selection below.
If no mark remains, the desired bound is immediate.

Choose a maximal direct-sum family $A$ of nonzero local spaces.
Put $t=|A|\ge1$ and $b=|\delta A|\le dt$.
The events $I_a$ for $a\in A$ are independent, so
\[
 \frac25\le\prod_{a\in A}\left(1-\frac{\PP(I_a)}2\right)
 \le(1-2^{-d})^t<e^{-t2^{-d}}.
\]
In particular, $t<2^d$.
If all remaining marks lie in $A$, the bound follows.
Otherwise $A$ is a nonempty proper vertex set of $Q$.
Every remaining mark outside $A$ is relevant in Lemma~\ref{lem:strings}, by maximality.
There are at most $2b$ exceptions and at most $2b$ strings.

Suppose a string has at least $(h_*+2)(s+3)$ marks.
Split these marks into $s+3$ consecutive groups of $h_*+2$ marks.
For each group take the region between its first and last marks from Lemma~\ref{lem:strings}.
Lift this region to the union of the sets contracted to its vertices.
The lifted regions are connected and pairwise disjoint.
Each contains $h_*$ disjoint cyclic sets $X_i$, so has rank at least $h_*$.
Its boundary has at most $2d\le4D$ edges.
At most $s$ regions meet $S$, and at most one meets $P$, since a nonempty connected $P$
lies in one contracted set. Thus two regions remain inside $G-(P\cup S)$.
Proposition~\ref{prop:blocks}(a) gives $q(G,P\cup S)\le\frac9{25}<\frac25$, a contradiction.

Count the omitted mark, the selected marks, the exceptions and the strings:
\[
 n\le1+t+2b+2b(h_*+2)(s+3).
\]
Here $b\le dt$, $t<2^{d}$, $s+3\le3(s+1)$ and $h_*+2=258$. So
$n\le2^{d}+2d\,2^{d}+1548\,d(s+1)2^{d}\le1551\,d(s+1)2^{d}$.
\end{proof}

\paragraph{Raising the girth.}

\begin{theorem}[Raising the girth]\label{thm:raise}
Let $G$ be a connected cubic multigraph and let $P$ be empty or induce a connected subgraph.
Suppose $q(G,P)\ge\frac25$ and $G-P$ has girth at least $D_*$.
Put $D_j=2^jD_*$ for $j\ge0$.
There are nested sets $S_0\subseteq S_1\subseteq\cdots$ outside $P$ such that
\begin{enumerate}
\item[\textup{(i)}] $S_j$ is the union of a maximal family of pairwise disjoint cycles of length
at most $D_j$ in $G-P$; hence $G-(P\cup S_j)$ has girth greater than $D_j$;
\item[\textup{(ii)}] $\log(|S_j|+1)\le2D_j+6\log^{2}D_j$.
\end{enumerate}
\end{theorem}

\begin{proof}
Put $S_{-1}=\emptyset$.
Having chosen $S_{j-1}$, take a maximal disjoint family $\mathcal Q_j$ of cycles of length
at most $b=D_j$ in $G-(P\cup S_{j-1})$, and add its vertices to $S_{j-1}$ to form $S_j$.
The exterior girth is at least $a=D_*$ for $j=0$, and at least $a=D_{j-1}+1$ otherwise;
thus $b\le2a$.
Each cycle $C$ has $|\delta V(C)|\le|C|\le b$.
Also $q(G,P\cup S_{j-1})\ge q(G,P)\ge\frac25$.
Proposition~\ref{prop:packing} gives
\[
 |\mathcal Q_j|\le1551\,b(|S_{j-1}|+1)2^b,
 \qquad |S_j|+1\le1552\,b^22^b(|S_{j-1}|+1).
\]
The union of all families chosen so far is maximal at scale $D_j$, which proves~(i).
For~(ii), take logarithms in the recursion. As $\log1552<11$, the step $i$ adds at most
$D_i+2\log D_i+11$. There are $j+1\le\log D_j$ steps, and $\sum_{i\le j}D_i<2D_j$. So
\[
 \log(|S_j|+1)\le2D_j+\log D_j\,\bigl(2\log D_j+11\bigr)\le2D_j+6\log^{2}D_j,
\]
because $\log D_j\ge\log7>\frac{11}4$.
\end{proof}

\subsection{Links, a second moment and a chain}\label{sec:chain}

\paragraph{Components between two connected sets.}
Lemma~\ref{lem:smallcuts} bounds the number of disjoint cyclic regions with small boundaries, and
it needs an excess above $\half$. For components that join the same two connected sets there is a
bound without an excess.

\begin{lemma}[Parallel links]\label{lem:parallel}
Let $\Gamma$ be a connected multigraph, let $A$ and $B$ be disjoint regions, and let $P$ be a set of
vertices with $q(\Gamma,P)\ge\frac25$. Let $t\ge0$ be an integer. Then at most $5\cdot2^{t}$
components $Y$ of $\Gamma-(A\cup B)$ have all of the following properties: $Y$ has an edge to $A$
and an edge to $B$; $Y$ avoids $P$; $Y$ contains a cycle; $|\delta Y|\le t$.
\end{lemma}

\begin{proof}
Let $Y_0,Y_1,\dots,Y_n$ be such components. For $1\le i<j\le n$ the graph $\Gamma-Y_i-Y_j$ is
connected. Indeed, it contains the connected set $A\cup B\cup Y_0$, and every other component of
$\Gamma-(A\cup B)$ is connected and has an edge to $A\cup B$. No edge joins $Y_i$ and $Y_j$. So the
regions $Y_i$ and $Y_j$ have exactly one link, and by Lemma~\ref{lem:link}
\[
  \PP(\cO_{Y_i}\cap\cO_{Y_j})=\PP(\cO_{Y_i})\,\PP(\cO_{Y_j}).
\]
Let $E_i$ be the event that $\cO_{Y_i}$ holds and $\xi\cap E(\Gamma[Y_i])$ is not empty. By the last
part of Lemma~\ref{lem:link}, the events $E_i$ and $E_j$ are independent as well, and
$\PP(E_i)\ge2^{-t}\cdot\half$. If $n=0$, the bound holds. Otherwise let $N$ be the number of the
events $E_1,\dots,E_n$ that occur, and put $\nu=\EE N\ge n\,2^{-t-1}>0$. Then
$\EE N^{2}\le\nu^{2}+\nu$, and the Cauchy--Schwarz inequality gives
\[
  \PP(N=0)\;\le\;1-\frac{\nu^{2}}{\EE N^{2}}\;\le\;\frac1{1+\nu}.
\]
If $\xi\cap E(\Gamma-P)$ is a forest, then no $E_i$ occurs, because $Y_i$ avoids $P$ and a nonempty
even subgraph is not a forest. So $\frac25\le q(\Gamma,P)\le1/(1+\nu)$. Hence $\nu\le\frac32$ and
$n\le3\cdot2^{t}$. With $Y_0$ there are at most $3\cdot2^{t}+1$ components.
\end{proof}

\begin{lemma}[Few links]\label{lem:linkshalf}
Let $G$ be a connected cubic multigraph, and let $P$ be a nonempty set of vertices such that $G[P]$
is connected and $q(G,P)\ge\frac25$. Let $S$ be the union of the vertex sets of a maximal family of
pairwise disjoint cycles of length at most $D$ in $G-P$, and put $J=G-(P\cup S)$. Let
$8\le D\le L$, and let $t$ be an integer with $1\le t\le D/4$. Then any two disjoint cycles $C,C'$
of $J$ of length at most $L$ satisfy
\[
  \ell\bigl(V(C),V(C')\bigr)\;\le\;1+\frac{2L}{t+1}+5\cdot2^{t}+289\Bigl(\frac LD\Bigr)^{2}.
\]
\end{lemma}

\begin{proof}
By the maximality of the family, $J$ is a simple graph of girth greater than $D$. So we are in the
setting of Section~\ref{sec:moderate}, with $\Gamma=G$ and $U=P\cup S$. Put
$U'=V(C)\cup V(C')$. Every vertex of $U'$ lies on exactly one edge outside its own cycle, so at
most $2L$ edges leave $U'$. At most one component of $G-U'$ meets $P$, because $G[P]$ is connected
and disjoint from $U'$. At most $2L/(t+1)$ components have more than $t$ edges to $U'$. Of the
other components that are joined to both cycles, at most $5\cdot2^{t}$ contain a cycle, by
Lemma~\ref{lem:parallel} with $A=V(C)$ and $B=V(C')$. The remaining components that are joined to
both cycles induce trees, avoid $P$ and have at most $t$ edges to $U'$. They also avoid $S$:
each cycle of the family is disjoint from $U'$, so it lies in one component of $G-U'$, and this
component contains a cycle. By Lemma~\ref{lem:treelinks}, at most $289(L/D)^{2}$ links are such
components or edges.
\end{proof}

\paragraph{Pairs with exactly two links.}
Let $G$, $P$, $S$, $J$, $D$ and $L$ be as in Lemma~\ref{lem:linkshalf}. As in its proof we are in
the setting of Section~\ref{sec:moderate}. So $s=\lfloor(D-1)/2\rfloor$, the set $\mathcal{C}$
consists of the cycles of $J$ of length at most $L$, and the numbers $a_C$, $\mu$ and $W_2$ and the
random variables $X_C$ and $Z$ are defined there.

\begin{lemma}[Two links]\label{lem:twolinks}
Suppose that $\ell(V(C),V(C'))\le b$ for all disjoint $C,C'\in\mathcal{C}$, and put
$\eta=2^{b}L^{3}2^{-s}$. Then
\[
  W_2\;\ge\;\tfrac23\,\mu^{2}-(1+\eta)\,\mu .
\]
\end{lemma}

\begin{proof}
By Lemma~\ref{lem:moment}, $\EE Z^{2}\le\mu^{2}+W_2+(1+\eta)\mu$. If $\xi\cap E(G-P)$ is a forest,
then $\xi$ contains no cycle of $J$, so $Z=0$. Hence $\PP(Z\ne0)\le\frac35$. The Cauchy--Schwarz
inequality $\mu^{2}\le\EE Z^{2}\cdot\PP(Z\ne0)$ gives $\EE Z^{2}\ge\frac53\mu^{2}$.
\end{proof}

\begin{idea}{from pairs with two links to a chain}
By Lemma~\ref{lem:twolinks} there are many pairs of cycles with exactly two links. Let $C$ and
$C'$ be such a pair. One of the two links avoids $p$, and then $C'$ separates a neighbour $u$ of
$C$ from $p$. Many disjoint cycles $C'$ have this property for the same $C$ and $u$. They are
nested (Figure~\ref{fig:chain}), and the parts between them are the pieces that we need. The
argument requires that $C'$ lies in the component of $p$ in $G-V(C)$. By symmetry, at least half
of the weight of the pairs with two links has this property.
\end{idea}

\paragraph{Separators are nested.}

\begin{lemma}[Nested separators]\label{lem:nested}
Let $u$ and $p$ be two vertices of a connected graph $H$. Let $C_1$ and $C_2$ be disjoint connected
sets of vertices that avoid $u$ and $p$, such that $u$ and $p$ lie in different components of
$H-C_i$ for $i=1,2$. Let $U_i$ be the vertex set of the component of $u$ in $H-C_i$. Then exactly
one of the following holds: $C_1\subseteq U_2$ and $U_1\subseteq U_2$; or $C_2\subseteq U_1$ and
$U_2\subseteq U_1$.
\end{lemma}

\begin{proof}
Take a path from $u$ to $p$, and let $z$ be its last vertex before the first vertex in $C_1$. The
part of the path up to $z$ lies in $U_1$. In the same way there is a path from $p$ to a vertex next
to $C_1$ that avoids $C_1$. The set $C_2$ is connected and disjoint from $C_1$, so it lies in one
component of $H-C_1$. If this component is neither $U_1$ nor the component of $p$, then the two
partial paths and a path inside $C_1$ form a walk from $u$ to $p$ that avoids $C_2$. This is
impossible. So $C_2$ lies in $U_1$ or in the component of $p$ in $H-C_1$, and the same holds with
the indices exchanged.

Suppose that $C_2$ lies in the component of $p$ in $H-C_1$. Then $C_2$ avoids $U_1$. So every
vertex of $U_1$ is joined to $u$ outside $C_2$, and so is every vertex of $C_1$, through $z$ and
the connected set $C_1$. Hence $U_1\cup C_1\subseteq U_2$. If $C_2$ lies in $U_1$ and $C_1$ lies in
$U_2$, then the path from $p$ to a vertex next to $C_1$, a path inside $C_1$, and a path inside
$U_2$ from $C_1$ to $u$ give a walk from $p$ to $u$ outside $C_2$. This is impossible. So if $C_2$
lies in $U_1$, then $C_1$ lies in the component of $p$ in $H-C_2$, and
$U_2\cup C_2\subseteq U_1$ by the first case with the indices exchanged. The two alternatives
exclude each other, because $C_1\subseteq U_2$ and $U_2\subseteq U_1$ would give
$C_1\subseteq U_1$.
\end{proof}

So a family of pairwise disjoint connected sets that separate $u$ from $p$ is linearly ordered:
write $C_1<C_2$ in the first case of the lemma. The order is transitive, because it is the order
of the sets $U_i$ by inclusion.

\begin{lemma}[Blocks]\label{lem:blocks}
In the situation of Lemma~\ref{lem:nested}, let $C_1<C_2$ and put
$X=U_2\sm(U_1\cup C_1)$. Then every edge of $H$ that leaves $X$ ends in $C_1$ or in $C_2$. If
$C_1<C<C_2$ for a third set $C$ of the family, then $C\subseteq X$. If $C_1<C_2\le C_3<C_4$, then
the sets $U_2\sm(U_1\cup C_1)$ and $U_4\sm(U_3\cup C_3)$ are disjoint.
\end{lemma}

\begin{proof}
An edge that leaves $U_2$ ends in $C_2$, because $U_2$ is a component of $H-C_2$. No edge joins
$X$ to $U_1$: its end in $X$ would lie in the component $U_1$ of $H-C_1$. So an edge from $X$ to
$U_1\cup C_1$ ends in $C_1$. For the second statement,
$C\subseteq U_2$, and $C$ avoids $U_1\cup C_1$ because $U_1\cup C_1\subseteq U_C$, the component of
$u$ in $H-C$. For the third, $U_2\subseteq U_3$.
\end{proof}

\paragraph{The chain.}

\begin{lemma}[Chain]\label{lem:chain}
In the situation of Lemma~\ref{lem:twolinks}, suppose that
\[
  \mu>30,\qquad\eta\le\frac12,\qquad L^{2}2^{-s}\le\frac12 .
\]
Fix a vertex $p\in P$. Then there are a cycle $C\in\mathcal C$, a vertex $u\notin V(C)$ that has a
neighbour on $C$, and pairwise disjoint cycles $C_1,\dots,C_T\in\mathcal C$ with
\[
  T\;\ge\;\frac{8\cdot2^{s}}{L^{3}},
\]
such that each $C_i$ is disjoint from $C$ and avoids $u$, and $u$ and $p$ lie in the same component
of $G-V(C)$ and in different components of $G-V(C)-V(C_i)$.
\end{lemma}

\begin{proof}
For disjoint $C,C'\in\mathcal C$ say that $C'$ is \emph{far behind} $C$ if $C'$ does not lie in the
component of $p$ in $G-V(C)$. This relation is antisymmetric. Let $C'$ be far behind $C$. A path
from $p$ to its first vertex on $C$ stays in the component of $p$ in $G-V(C)$ before it ends. So it
avoids $C'$. As $C$ is connected and disjoint from $C'$, all of $C$ lies in the component of $p$ in
$G-V(C')$. So $C$ is not far behind $C'$. Hence the ordered pairs $(C,C')$ with $C'$ far behind $C$
have total weight at most $W_2/2$ when restricted to the exactly-two-link pairs.
Indeed, the two-link relation is symmetric and opposite orientations have the same weight;
at most one orientation is far behind.

Fix $C$. Every vertex of $C$ has exactly one edge end outside $C$, so at most $L$ vertices outside
$C$ have a neighbour on $C$. A cycle $C'$ that is joined to $C$ by an edge passes through one of
them. By~\eqref{eq:load} these cycles have total weight at most $L^{2}2^{-s}\le\half$. So the
ordered pairs of disjoint cycles that are joined by an edge have total weight at most $\mu/2$.

Call an ordered pair $(C,C')$ \emph{good} if the cycles are disjoint, have exactly two links, are
not joined by an edge, and $C'$ is not far behind $C$. By Lemma~\ref{lem:twolinks} and the two
bounds above, the good pairs have total weight at least $W_2/2-\mu/2$, hence at least
\[
  \frac13\mu^2-\left(1+\frac\eta2\right)\mu
  \;>\;\left(10-\frac54\right)\mu\;>\;8\mu .
\]
So for some $C\in\mathcal C$ the cycles $C'$ with $(C,C')$ good have total weight more than $8$.
Fix this $C$, and put $H=G-V(C)$.

Let $(C,C')$ be good. The two links are components $Y_1,Y_2$ of $G-V(C)-V(C')$ that are joined to
both cycles, because no edge joins the cycles. The vertex $p$ lies in at most one of them. So one of
them, say $Y$, avoids $p$. Let $u$ be a vertex of $Y$ with a neighbour on $C$. The cycle $C'$ lies
in the component of $p$ in $H$, and $Y$ is connected and joined to $C'$. So $u$ lies in that
component as well. The set $Y$ is a component of $H-V(C')$ that contains $u$ and not $p$. So $C'$
separates $u$ from $p$ in $H$.

There are at most $L$ choices for $u$. So for some vertex $u$ the cycles $C'$ that separate $u$
from $p$ in $H$, as above, have total weight at least $8/L$. Take a maximal family
$C_1,\dots,C_T$ of pairwise disjoint cycles among them. Every cycle of the collection meets one of
the at most $TL$ vertices of the family. By~\eqref{eq:load},
\[
  \frac8L\;\le\;TL\cdot L\,2^{-s}.\qedhere
\]
\end{proof}

\subsection{Proof of Theorem~\ref{thm:half}}\label{sec:halfproof}

Let $D_*$ and $C_0$ be as in Lemma~\ref{lem:affine}. Let $G$ and $P$ be as in the theorem, and
suppose that $q(G,P)\ge\frac25$. We derive a contradiction when $r$ is large. All limits below are
for $r\to\infty$, and they are uniform in $G$ and $P$. Put
\[
  x=\log r,\qquad z=\log x,\qquad A=6\,e^{124},\qquad L=\lfloor Ax\rfloor,\qquad
  t=\lfloor z-2\log z\rfloor .
\]

\paragraph{Two girth scales.}
Let $D_j=2^{j}D_{*}$ and let $S_0\subseteq S_1\subseteq\cdots$ be the sets of
Theorem~\ref{thm:raise}. Let $D^{-}$ be the first $D_j$ that is at least $\sqrt x$, and let $D^{+}$
be the first $D_j$ that is at least $x/16$. For large $r$,
\[
  \sqrt x\le D^{-}<2\sqrt x,\qquad \frac x{16}\le D^{+}<\frac x8,\qquad D^{-}<D^{+}.
\]
Let $S^{-}\subseteq S^{+}$ be the corresponding sets, and put $J^{\pm}=G-(P\cup S^{\pm})$. Then
$J^{-}$ has girth greater than $D^{-}$, $J^{+}$ has girth greater than $D^{+}$, and
$J^{+}\subseteq J^{-}$. By Theorem~\ref{thm:raise}(ii),
\begin{equation}\label{eq:twosizes}
  |S^{-}|\;\le\;2^{4\sqrt x+6\log^{2}(2\sqrt x)}=r^{o(1)},\qquad
  |S^{+}|\;\le\;2^{x/4+6\log^{2}x}=r^{1/4+o(1)}.
\end{equation}
We use the small scale $D^{-}$ for the girth of the final pieces, and the large scale $D^{+}$
for a long chain.

\paragraph{Cycles at the large scale.}
We apply Section~\ref{sec:chain} with $S=S^{+}$, $J=J^{+}$, $D=D^{+}$ and
$s=\lfloor(D^{+}-1)/2\rfloor\ge x/32-1$. The hypotheses of Lemma~\ref{lem:linkshalf} hold for large
$r$: $8\le D^{+}\le L$ and $1\le t\le z\le D^{+}/4$. So every two disjoint cycles of $J^{+}$ of
length at most $L$ have at most
\[
  b\;=\;1+\frac{2L}{t+1}+5\cdot2^{t}+289\Bigl(\frac L{D^{+}}\Bigr)^{2}
  \;\le\;1+\frac{2Ax}{z-2\log z}+\frac{5x}{z^{2}}+289\,(16A)^{2}\;\le\;\frac{3Ax}{z}
\]
links. Hence
\[
  \eta=2^{b}L^{3}2^{-s}\;\le\;r^{3A/z}\,(Ax)^{3}\,2\,r^{-1/32}=r^{-1/32+o(1)}=o(1),
  \qquad L^{2}2^{-s}=o(1).
\]

The graph $J^{+}$ has $v=2r-2-|P|-|S^{+}|=2r-o(r)$ vertices, and its deficit is
$d_J\le3(|P|+|S^{+}|)=r^{1/4+o(1)}$. Let $\mu$ be the total weight of its cycles of length at most
$L$. By Lemma~\ref{lem:weight} with $\ell_0=\lceil6x\rceil$,
\[
  \Bigl(1+\tfrac32\,L^{2}\,2^{-s}\Bigr)\,\mu\;\ge\;
  \Bigl(1-\frac{L\,d_J}{2v}-3v\,r^{-3}\Bigr)\sum_{\ell}\frac1{2\ell},
\]
where the sum is over the even $\ell$ with $6x\le\ell\le L$. This sum is
$(\frac14-o(1))\ln(A/6)=31-o(1)$, and the two factors in brackets are $1+o(1)$. Thus $\mu>30$ for
large $r$.

\paragraph{The chain.}
Fix $p\in P$. For large $r$ the hypotheses of Lemma~\ref{lem:chain} hold. So there are a cycle $C$
of $J^{+}$ of length at most $L$, a vertex $u$, and pairwise disjoint cycles $C_1,\dots,C_T$ of
$J^{+}$ of length at most $L$ with
\[
  T\;\ge\;\frac{8\cdot2^{s}}{L^{3}}\;\ge\;\frac{4}{A^{3}}\cdot\frac{r^{1/32}}{x^{3}},
\]
such that each $C_i$ separates $u$ from $p$ in $G-V(C)$. Let $H$ be the component of $G-V(C)$
that contains $u$ and $p$. Each $V(C_i)$ is connected and separates $u$ from $p$, so it lies in $H$.
By Lemma~\ref{lem:nested} we may number the cycles so that $C_1<C_2<\dots<C_T$. Let $U_i$ be the
vertex set of the component of $u$ in $H-V(C_i)$.

\begin{figure}[ht]
\centering
\begin{tikzpicture}[x=1cm,y=1cm]
  % the cycle C and the vertex u
  \draw[hot] (0,0) ellipse [x radius=0.32, y radius=0.95];
  \node[note, text=warm] at (0,1.25) {$C$};
  \node[vtx] (u) at (0.85,0.25) {};
  \draw[edge] (0.31,0.25)--(u);
  \node[note, above=2pt] at (u) {$u$};
  % blocks
  \fill[accent!10] (2.4,-1.0) rectangle (5.2,1.0);
  \fill[accent!10] (5.8,-1.0) rectangle (8.6,1.0);
  \node[note, text=accent] at (3.8,-1.3) {the block $X_0$};
  \node[note, text=accent] at (7.2,-1.3) {the block $X_1$};
  % separating cycles
  \foreach \x in {2.1,5.5,8.9} \draw[cool] (\x,0) ellipse [x radius=0.26, y radius=0.95];
  \foreach \x in {2.9,3.5,4.1,4.7,6.3,6.9,7.5,8.1,9.7,10.3}
    \draw[accent, line width=0.45pt] (\x,0) ellipse [x radius=0.2, y radius=0.85];
  \node[note, text=accent] at (2.1,1.25) {$C_1$};
  \node[note, text=accent] at (5.5,1.25) {$C_{w+1}$};
  \node[note, text=accent] at (8.9,1.25) {$C_{2w+1}$};
  \node[note] at (10.95,0) {$\cdots$};
  % the marked set P
  \draw[draw=slate!70, fill=black!4, line width=0.6pt] (12.2,0) ellipse [x radius=0.55, y radius=0.45];
  \node[text=slate] at (12.2,0) {\small$P$};
  \node[vtx] (pp) at (11.92,0.05) {};
  \node[note, above=7pt] at (12.2,0.3) {$p\in P$};
  \draw[slate!60, line width=0.4pt, dash pattern=on 2pt off 2pt] (u)--(1.8,0.25) (11.3,0.05)--(pp);
\end{tikzpicture}
\caption{The chain of Lemma~\ref{lem:chain}. Each cycle $C_i$ separates $u$ from $p$ in $G-V(C)$, so
the cycles are nested. Every $w$-th cycle is the end of a block. A block has many vertices and at
most $2L$ boundary edges.}
\label{fig:chain}
\end{figure}

\paragraph{Blocks.}
Put $w=\lceil128A^{2}x\rceil+2$. For $l=0,1,\dots$ with $(l+1)w\le T-1$ let
\[
  X_l\;=\;U_{(l+1)w+1}\sm\bigl(U_{lw+1}\cup V(C_{lw+1})\bigr).
\]
There are at least $\lfloor(T-1)/w\rfloor\ge r^{1/32}/x^{5}$ such sets, for large $r$. By
Lemma~\ref{lem:blocks} they are pairwise disjoint. Each contains the $w-1$ cycles between its two
end cycles. These cycles are pairwise disjoint and each has more than $D^{+}$ vertices, so
\[
  |X_l|\;\ge\;(w-1)(D^{+}+1)\;\ge\;128A^{2}x\cdot\frac x{16}=8A^{2}x^{2}\;\ge\;8L^{2}.
\]
Every edge of $H$ that leaves $X_l$ ends on one of its two end cycles. An end cycle has at most $L$
vertices, and each of them has one edge end outside the cycle. So at most $2L$ edges of $H$ leave
$X_l$.

\paragraph{Clean blocks.}
Discard a set $X_l$ if it contains a vertex with a neighbour on $C$, or a vertex of $S^{-}$. At most
$L$ vertices outside $C$ have a neighbour on $C$, so this discards at most
$L+|S^{-}|=r^{o(1)}$ sets. For large $r$ at least two sets remain. Let $X$ be one of them.
\begin{itemize}
\item $X$ avoids $P$. The set $P$ is connected and disjoint from the cycles, and it contains $p$,
which is not in $U_i$ for any $i$. A path in $G[P]$ from $p$ to a vertex of $U_i$ would have to pass
through $C_i$ or through $C$.
\item No edge of $G$ joins $X$ to $V(C)$, and $H$ is a component of $G-V(C)$. So the edges of $G$
that leave $X$ are edges of $H$, and $|\delta X|\le2L$.
\item $X$ avoids $P\cup S^{-}$, so $G[X]$ is a subgraph of $J^{-}$, and its girth is greater than
$D^{-}$.
\end{itemize}

\paragraph{The contradiction.}
Two of the remaining sets satisfy the hypotheses of Proposition~\ref{prop:blocks}(b) with
$\Lambda=L$ and $D=D^{-}$: they are disjoint, they avoid $P$, each has at least $8L^{2}$ vertices
and at most $2L$ boundary edges, and $D^{-}\ge\sqrt x\ge C_0\log(L+2)$ for large $r$. Thus
$q(G,P)\le\frac9{25}<\frac25$. This contradiction proves Theorem~\ref{thm:half}. \qed

\section{Proof of Theorem~\ref{thm:exact}}\label{sec:proof}

Let $D_*$ be as in Theorem~\ref{thm:half} and $K_0$ as in Section~\ref{sec:poles}.
By~\eqref{eq:H},
\begin{equation}\label{eq:Hlower}
  H(r)\;\ge\;2^{r-1}+(r-2)\,2^{r/2-1}\qquad\text{for all }r\ge3 .
\end{equation}
Choose $R\ge8$ so large that the following holds for every $r\ge R$:
Proposition~\ref{prop:threecut}, Theorem~\ref{thm:fourcut}, Theorem~\ref{thm:girth} with
$g_0=D_*-1$, and Theorem~\ref{thm:half} apply to the rank $r$; and $r-2>10K_0^{2}$ and
$2\log(2r)+2\le3\log r$.

Let $r\ge R$. A ring $R_r$ has $H(r)$ cycles by Corollary~\ref{cor:ringmax}, so $\psi(r)\ge H(r)$.
By Lemma~\ref{lem:cubic}, some connected cubic multigraph $G$ of rank $r$ has $\psi(r)$ cycles. We
show that $G$ is a ring or that $c(G)\le2^{r-1}+5K_0^{2}\,2^{r/2}$. In the first case
$c(G)\le H(r)$ by Corollary~\ref{cor:ringmax}. In the second case $c(G)<H(r)$ by~\eqref{eq:Hlower},
because $r-2>10K_0^{2}$. So $\psi(r)\le H(r)$.

\step{Case 1: $G$ has a loop $e$.} Every cycle of $G$ other than $e$ is a cycle of $G-e$, which has
rank $r-1$. So $c(G)\le1+(2^{r-1}-1)$.

\step{Case 2: $G$ has no loop, and it has a cyclic cut of size at most three.} Then
$c(G)\le2^{r-1}+5K_0^{2}\,2^{r/2}$ by Proposition~\ref{prop:threecut}.

\step{Case 3: $G$ has no loop, it has a cyclic cut of size four, and it has no smaller cyclic
cut.} Then $G$ is a ring or $c(G)<2^{r-1}$, by Theorem~\ref{thm:fourcut}.

In the other cases $G$ has no loop and no cyclic cut of size at most four. Let $g$ be its girth.

\step{Case 4: $g<D_*$.} Then $c(G)<2^{r-1}$ by Theorem~\ref{thm:girth}.

\step{Case 5: $g\ge D_*$.} Let $W$ be a shortest cycle, and put $P=V(W)$. As in the proof of
Corollary~\ref{cor:cycles}, $|P|=g\le2\log(2r)+2\le3\log r$. The graph $G[P]$ is connected, and
$G-P$ has girth at least $D_*$. So $q(G,P)<\frac25$ by Theorem~\ref{thm:half}.

Let $Z$ be a cycle of $G$ that meets $P$. Its edges in $G-P$ form a proper subset of $Z$, hence a
forest. So at most $q(G,P)\,2^{r}<\frac25\,2^{r}$ cycles meet $P$. The other cycles are cycles of
$G-P$, and there are at most $2^{\beta(G-P)}$ of them. The cycle $W$ has no chord. So $G-P$ arises
from $G$ by the deletion of $g$ vertices and $2g$ edges, and
\[
  \beta(G-P)=r-1-g+\comp(G-P).
\]
The graph $G$ has no cut of size one, because such a cut is cyclic (Section~\ref{sec:cuts}). So
every component of $G-P$ contains the ends of at least two of the $g$ edges that leave $P$. Hence
$\comp(G-P)\le g/2$ and $\beta(G-P)\le r-1-g/2\le r-4$, because $g\ge7$. Therefore
$c(G)<\bigl(\frac25+\frac1{16}\bigr)2^{r}<2^{r-1}$. \qed


\begin{thebibliography}{99}\small

\bibitem{AACST} A.~Alahmadi, R.~E.~L.~Aldred, R.~dela~Cruz, P.~Sol\'e and C.~Thomassen, The maximum
number of minimal codewords in long codes, \emph{Discrete Appl.\ Math.}\ \textbf{161} (2013),
424--429. arXiv:1201.6043.

\bibitem{AT} R.~E.~L.~Aldred and C.~Thomassen, On the maximum number of cycles in a planar graph,
\emph{J.~Graph Theory} \textbf{57} (2008), 255--264.

\bibitem{ES} R.~C.~Entringer and P.~J.~Slater, On the maximum number of cycles in a graph,
\emph{Ars Combin.}\ \textbf{11} (1981), 289--294. A scan is linked from the entry A006188 of the
OEIS, \url{https://oeis.org/A006188}.

\bibitem{Forest} KNT, Astra and Fable~5.2, The forest proportion of a random even subgraph,
companion manuscript, November 2026.

\bibitem{HK} P.~Horak and D.~Kim, Connected cubic graphs with the maximum number of perfect
matchings, \emph{J.~Graph Theory} \textbf{99} (2022), 671--690. arXiv:2006.13459.

\bibitem{KNT} KNT, The minimum number of edges in a pancyclic graph, manuscript, September 2026,
\url{https://jehlp.net/erdos1016/}.

\bibitem{Repo} KNT, Astra and Fable~5.2, Lean~4 formalisation of the forest estimate, of the
results of~\cite{KNT} and of Theorem~\ref{thm:exact}, \url{https://github.com/JWKNT/erdos1016}.

\end{thebibliography}

\end{document}
